\newcommand{\langue}{french}
\newcommand{\nomchapitre}{Chapitre}
\newcommand{\titrecourt}{MANUEL DE RÉFÉRENCE}
\input{version.tex}
% ================================================================
%  Préambule du manuel de Coupole (FR et EN), d'après le format
%  maison : Latin Modern, titres sans-serif colorés, encadrés
%  tcolorbox, fontawesome5.  \langue{french} ou \langue{english}
%  est défini par le fichier principal avant \input.
% ================================================================
\documentclass[11pt,a4paper]{article}

\usepackage[T1]{fontenc}
\usepackage[utf8]{inputenc}
\usepackage[\langue]{babel}
\usepackage{lmodern}
\usepackage{textcomp}
\usepackage{amsmath,amssymb}
\usepackage{mathtools}
\usepackage{siunitx}
\sisetup{group-minimum-digits=4,separate-uncertainty=true,per-mode=symbol}
\usepackage{cancel}
\usepackage{microtype}
\DeclareUnicodeCharacter{2248}{\ensuremath{\approx}}
\DeclareUnicodeCharacter{2212}{\ensuremath{-}}
\DeclareUnicodeCharacter{00D7}{\ensuremath{\times}}
\DeclareUnicodeCharacter{2192}{\ensuremath{\rightarrow}}
\DeclareUnicodeCharacter{00A0}{~}
\DeclareUnicodeCharacter{202F}{\,}
\DeclareUnicodeCharacter{2033}{\ensuremath{''}}
\DeclareUnicodeCharacter{2032}{\ensuremath{'}}
\DeclareUnicodeCharacter{2264}{\ensuremath{\le}}
\DeclareUnicodeCharacter{2265}{\ensuremath{\ge}}
\DeclareUnicodeCharacter{2194}{\ensuremath{\leftrightarrow}}
\DeclareUnicodeCharacter{00B2}{\textsuperscript{2}}
\DeclareUnicodeCharacter{2026}{\ldots}
\DeclareUnicodeCharacter{00B0}{\textdegree}

\usepackage[a4paper,margin=2cm,headheight=24pt]{geometry}
\usepackage{graphicx}
\graphicspath{{img/}}
\usepackage[table]{xcolor}
\usepackage[most]{tcolorbox}
\usepackage{booktabs}
\usepackage{tabularx}
\usepackage{array}
\usepackage{multicol}
\usepackage{enumitem}
\usepackage[nobottomtitles*]{titlesec}
\usepackage{needspace}
\usepackage{fancyhdr}
\usepackage{fontawesome5}
\usepackage{tikz}
\usetikzlibrary{arrows.meta,calc,decorations.pathreplacing,patterns,angles,quotes}
\usepackage{pgfplots}
\pgfplotsset{compat=1.18}
\usepgfplotslibrary{fillbetween}
\usepackage{pdfpages}
\usepackage[colorlinks=true,linkcolor=sky,urlcolor=sky,bookmarksopen=true,bookmarksopenlevel=1,
  pdftitle={Coupole},
  pdfauthor={ARP273-ROSE}]{hyperref}
\usepackage{bookmark}
% Les notes posées dans des tableaux (annexes) par \footnotemark/\footnotetext cassent les ancres
% de hyperref : les appels de note ne sont donc pas des liens.
\hypersetup{hyperfootnotes=false}
\usepackage{xurl}
\usepackage{fvextra}
\fvset{breaklines=true,fontsize=\small}

% ---------------------------------------------------------------- Langue des titres d'encadrés
\newcommand{\tr}[2]{\ifdefstring{\langue}{english}{#2}{#1}}

% ---------------------------------------------------------------- Couleurs
\definecolor{sky}{HTML}{1C6DB4}
\definecolor{skyl}{HTML}{E7F0F9}
\definecolor{alert}{HTML}{B5382B}
\definecolor{alertl}{HTML}{FBEAE7}
\definecolor{night}{HTML}{1B2A4A}
\definecolor{ink}{HTML}{1F2430}
\definecolor{summit}{HTML}{6B7A8F}
\definecolor{gold}{HTML}{B8860B}
\definecolor{goldl}{HTML}{FBF3DC}
\definecolor{pine}{HTML}{2C7A55}
\definecolor{pinel}{HTML}{E2F1E9}
\definecolor{prune}{HTML}{5B3F8C}
\definecolor{prunel}{HTML}{EEE8F6}
\AtBeginDocument{\color{ink}}

% ---------------------------------------------------------------- Titres
% \section = chapitre du cours (« Chapitre 1 : ... »), \subsection = I, II, III..., \subsubsection = A, B, C...
\titleformat{\section}[display]
  {\sffamily\color{night}}
  {\large\bfseries\color{alert}\MakeUppercase{\sectionname}~\thesection}{2pt}{\LARGE\bfseries}
  [{\vspace{2pt}\color{night!35}\titlerule[1.3pt]}]
\titlespacing*{\section}{0pt}{6pt}{10pt}
\newcommand{\sectionname}{\nomchapitre}
\renewcommand{\thesubsection}{\arabic{section}.\arabic{subsection}}
\renewcommand{\thesubsubsection}{}
\titleformat{\subsection}
  {\large\bfseries\sffamily\color{sky}}{\thesubsection}{0.6em}{}
\titlespacing*{\subsection}{0pt}{12pt}{5pt}
\titleformat{\subsubsection}
  {\normalsize\bfseries\sffamily\color{alert}}{}{0em}{}
\titlespacing*{\subsubsection}{0pt}{9pt}{3pt}
\setcounter{secnumdepth}{3}
\setcounter{tocdepth}{2}

\usepackage{etoolbox}
\pretocmd{\section}{\clearpage}{}{}
\pretocmd{\subsection}{\needspace{8\baselineskip}}{}{}
\pretocmd{\subsubsection}{\needspace{5\baselineskip}}{}{}
\renewcommand{\bottomtitlespace}{0.15\textheight}
\widowpenalty=10000
\clubpenalty=10000
\raggedbottom

% ---------------------------------------------------------------- En-têtes
\pagestyle{fancy}
\fancyhf{}
\renewcommand{\headrulewidth}{0.4pt}
\renewcommand{\headrule}{\hbox to\headwidth{\color{night!40}\leaders\hrule height \headrulewidth\hfill}}
\fancyhead[L]{\sffamily\small\color{night}\textbf{COUPOLE · \titrecourt}}
\fancyhead[R]{\sffamily\small\color{summit}\partiecourante}
\fancyfoot[L]{\sffamily\scriptsize\color{summit}Coupole \versioncoupole{} · GPL-3.0 · ARP273-ROSE}
\fancyfoot[R]{\sffamily\small\color{summit}\thepage}
\newcommand{\partiecourante}{}
\newcommand{\partie}[1]{\clearpage\renewcommand{\partiecourante}{#1}}


% ---------------------------------------------------------------- Parties
\titleformat{\part}[display]{\centering\sffamily\color{night}}
  {\Large\bfseries\color{alert}PARTIE \thepart}{8pt}{\Huge\bfseries}
\pretocmd{\part}{\clearpage}{}{}
\setcounter{tocdepth}{2}

% ---------------------------------------------------------------- Encadrés du site (une icône par type de section)
\tcbset{maison/.style={enhanced,breakable,boxrule=0.9pt,arc=2.5pt,left=2.5mm,right=2.5mm,top=1.5mm,bottom=1.5mm,
  fonttitle=\bfseries\sffamily\small,coltitle=white,
  before upper={\widowpenalties 4 10000 10000 10000 0 \clubpenalties 4 10000 10000 10000 0}}}
% \typesite{env}{icône}{libellé}{fond}{cadre}
\newcommand{\typesite}[5]{\newtcolorbox{#1}[1][]{maison,colback=#4,colframe=#5,title={#2~~#3\ifx\relax##1\relax\else~·~##1\fi}}%
  \BeforeBeginEnvironment{#1}{\par\needspace{6\baselineskip}}}
\typesite{cdefinition}{\faBook}{\tr{Définition}{Definition}}{skyl}{sky}
\typesite{crappel}{\faUndo}{\tr{Rappel}{Reminder}}{goldl}{gold!85!black}
\typesite{cremarque}{\faInfoCircle}{\tr{Remarque}{Note}}{skyl}{sky!70!black}
\typesite{cattention}{\faExclamationTriangle}{\tr{Attention}{Warning}}{alertl}{alert}
\typesite{cexemple}{\faLightbulb}{\tr{Exemple}{Example}}{pinel}{pine}
\typesite{cexercice}{\faPencil*}{\tr{Exercice}{Exercise}}{white}{pine}
\typesite{cdemonstration}{\faCalculator}{\tr{Démonstration}{Proof}}{prunel}{prune}
\typesite{cconseil}{\faThumbsUp}{\tr{Conseil}{Tip}}{pinel}{pine!80!black}
\typesite{ccomplement}{\faPlusCircle}{\tr{Complément}{Supplement}}{prunel}{prune!80!black}
\typesite{censavoirplus}{\faSearchPlus}{\tr{En savoir plus}{Further reading}}{prunel}{prune!70!black}
\typesite{cactivite}{\faFlask}{\tr{Activité}{Activity}}{goldl}{gold}
\typesite{cqcm}{\faCheckSquare}{\tr{QCM}{Quiz}}{white}{sky}
\typesite{cobjectifs}{\faBullseye}{\tr{Objectifs}{Goals}}{skyl}{night}
\typesite{cprerequis}{\faListUl}{\tr{Prérequis}{Prerequisites}}{skyl}{night}
\typesite{ccadre}{\faStar}{\tr{À retenir}{Key point}}{goldl}{gold}
\newtcolorbox{ccode}{maison,colback=night!4,colframe=night!50,fontupper=\small\ttfamily,title={\faTerminal~~\tr{Commande}{Command}}}

\newcommand{\intertitre}[2][]{\par\needspace{4\baselineskip}\medskip\noindent{\sffamily\bfseries\color{night}#2}\par\smallskip}
\newcommand{\question}[1]{\par\smallskip\noindent{\sffamily\bfseries\color{pine}#1}~}
\newcommand{\auteurs}[1]{\par\noindent{\sffamily\small\color{summit}\faUserEdit~#1}\par\medskip}
\newcommand{\appli}{\textcolor{summit}{\faMousePointer~\textit{(application interactive en ligne)}}}
% \figsite{titre}{image}{légende}{crédit}
\newcommand{\figsite}[4]{\par\needspace{8\baselineskip}\begin{center}\begin{minipage}{\linewidth}\centering
  \if\relax\detokenize{#1}\relax\else{\sffamily\small\bfseries\color{sky}#1}\par\smallskip\fi
  \if\relax\detokenize{#2}\relax{\sffamily\small\color{summit}\faFilm~\textit{animation ou vidéo en ligne, non reproduite}}\par\else#2\par\fi
  \if\relax\detokenize{#3}\relax\else\smallskip{\small #3}\par\fi
  \if\relax\detokenize{#4}\relax\else{\scriptsize\color{summit}Crédit : #4}\par\fi
  \end{minipage}\end{center}}
\DeclareUnicodeCharacter{2009}{\,}
\DeclareUnicodeCharacter{2013}{--}
\DeclareUnicodeCharacter{2014}{---}
\DeclareUnicodeCharacter{2019}{'}
\DeclareUnicodeCharacter{2018}{`}
\DeclareUnicodeCharacter{201C}{``}
\DeclareUnicodeCharacter{201D}{''}
% \ressource{titre}{niveau · format · volume}{url}{description} — guide du portail
\newcommand{\ressource}[4]{\par\needspace{5\baselineskip}\noindent{\sffamily\bfseries\color{night}#1}\\
  {\sffamily\footnotesize\color{summit}#2 · \href{#3}{\faExternalLink*~lien}}\\
  {\small #4}\par\smallskip}

\renewcommand{\partiecourante}{Coupole \versioncoupole}
\newcommand{\capture}[2]{\par\needspace{10\baselineskip}\begin{center}\includegraphics[width=#2\linewidth]{#1}\end{center}}
\newcommand{\cmd}[1]{\texttt{#1}}

\begin{document}
% ---------------------------------------------------------------- page de garde
\begin{titlepage}
\centering
\vspace*{2.5cm}
\includegraphics[width=4.2cm]{../../logo/coupole_512.png}\par\vspace{1.2cm}
{\sffamily\bfseries\color{night}\fontsize{38}{44}\selectfont Coupole\par}
\vspace{0.4cm}
{\sffamily\Large\color{sky}Manuel de référence\par}
\vspace{0.3cm}
{\sffamily\large\color{summit}Version \versioncoupole{} --- \datecoupole\par}
\vspace{1.4cm}
\begin{minipage}{0.78\linewidth}\centering\sffamily\color{ink}
Boîte à outils libre pour les étudiants du DU « Explorer et Comprendre l'Univers » (Observatoire de Paris).\\[4pt]
Banque d'images de l'OHP · qualité des images · spectres et séries · sites et heures · cosmologie.
\end{minipage}
\vfill
{\sffamily\small\color{summit}Auteur : ARP273-ROSE · Licence GNU GPL version 3 ou ultérieure\\
Données : banque « OHP student observations », Observatoire de Paris / PADC (licence Etalab 2.0)\par}
\end{titlepage}

\tableofcontents

% =====================================================================
\section{Présentation}
Coupole est une application libre (GPL-3.0), écrite en Python 3.10+, qui fonctionne sous Windows, macOS et Linux,
avec une interface graphique (PyQt6) et une ligne de commande complète. Tout l'affichage existe en français et en
anglais ; la langue du système est détectée, et peut être imposée (menu \emph{Langue}, Préférences, ou
\cmd{--lang fr|en}).

\begin{center}\small
\begin{tabularx}{\linewidth}{@{}lX@{}}
\toprule
\textbf{Module} & \textbf{Fonction}\\
\midrule
Banque OHP & Inventaire, catalogue, téléchargement, contrôle astrométrique, en-têtes corrigés, conversion, lots empilables, anomalies, carte du ciel.\\
Qualité des images & Facultatif. FWHM, ellipticité, fond, bruit, gradient, saturation, traînées, échantillonnage.\\
Spectres et séries & Lecture et tracé de données 1D (spectres radio, courbes de lumière, tables).\\
Cosmologie & Du redshift aux distances, âges, volume, module de distance, échelle (Planck 2018 et autres).\\
Sites et heures & Base des sites, carte OpenStreetMap, heure UTC et heure locale.\\
Ma machine & Diagnostic matériel ; exemple minimal de module.\\
\bottomrule
\end{tabularx}
\end{center}

Principes : tout travail long s'exécute hors du fil graphique (l'interface reste réactive, avec progression et
annulation) ; un traitement interrompu reprend où il s'est arrêté ; rien n'est supprimé sans trace ; aucune
donnée n'est envoyée sans accord ; les composants facultatifs (ASTAP, SEP, CuPy) peuvent manquer.

% =====================================================================
\section{Installation}
\subsection{Paquets tout compris}
Les paquets embarquent Python et toutes les bibliothèques (PyQt6, numpy, astropy, zstandard, lz4, SEP) : rien
d'autre n'est à installer. Ils se mettent à jour d'eux-mêmes. ASTAP reste facultatif et n'est pas inclus.

\begin{center}\small
\begin{tabularx}{\linewidth}{@{}l>{\ttfamily\small}lX@{}}
\toprule
\textbf{Système} & \textbf{\normalfont Fichier} & \textbf{Installation}\\
\midrule
Windows 10/11 64 bits & Coupole-Setup-x.y.z.exe & Double-clic. Avertissement SmartScreen au premier lancement (installeur non signé) : « Informations complémentaires », « Exécuter quand même ». Installation dans le profil, sans droit administrateur.\\
macOS Apple Silicon & Coupole-x.y.z-macos-arm64.dmg & Glisser Coupole dans Applications. Premier lancement : clic droit, « Ouvrir ».\\
macOS Intel & Coupole-x.y.z-macos-x86\_64.dmg & Idem.\\
Ubuntu / Debian x86\_64, ARM64 & coupole\_x.y.z\_amd64.deb & \texttt{sudo apt install ./coupole\_x.y.z\_amd64.deb} (\texttt{\_arm64.deb} sur ARM) : apt installe les bibliothèques système de Qt ; commande \texttt{coupole}, entrée de menu. Ubuntu 22.04 ou plus récent, Debian 12 ou plus récent. Désinstaller : \texttt{sudo apt remove coupole}.\\
Autre Linux x86\_64 / ARM64 & Coupole-x.y.z-linux-*.tar.gz & Décompresser, lancer \texttt{installer.sh} (copie dans \texttt{\textasciitilde/.local/share/coupole}, entrée de menu, commande \texttt{coupole} par un lien \texttt{\textasciitilde/.local/bin/coupole} ; si ce dossier n'est pas dans le \texttt{PATH}, l'installeur le dit et donne la ligne à ajouter). Le lanceur \texttt{Coupole.sh} suit les liens symboliques (lien de lien compris).\\
\bottomrule
\end{tabularx}
\end{center}

\subsection{Avec le Python de l'ordinateur}
Les scripts du dossier des sources cherchent un Python 3.10 ou plus récent, créent un environnement virtuel,
installent les dépendances et créent les lanceurs. Une seule question : Entrée pour confirmer (\cmd{--oui} /
\cmd{-Oui} pour l'omettre).
\begin{ccode}
sh install.sh                \# Linux, macOS\\
install.bat                  \# Windows (double-clic ; lance install.ps1)
\end{ccode}
Sans Python, ou sans le module \texttt{venv} (Debian : paquet \texttt{python3-venv}), le script le dit et indique
comment l'obtenir. Emplacements : Linux \texttt{\textasciitilde/.local/share/coupole/venv} et lanceurs dans
\texttt{\textasciitilde/.local/bin} ; macOS \texttt{\textasciitilde/Library/Application Support/Coupole/venv} et
\texttt{\textasciitilde/Applications/Coupole.app} ; Windows \texttt{\%LOCALAPPDATA\%\textbackslash
Programs\textbackslash CoupolePython} et raccourcis Bureau et menu Démarrer.

\subsection{pip et pipx}
\begin{ccode}
pipx install git+https://github.com/ARP273-ROSE/coupole\\
pip install .                 \# depuis les sources ; pip install ".[qualite]" ajoute SEP\\
python -m coupole             \# interface ; coupole --help : ligne de commande
\end{ccode}
Extras : \texttt{qualite} (SEP), \texttt{validation} (lxml, xisf : validation XSD et second lecteur),
\texttt{gpu} (CuPy), \texttt{test} (pytest).

\subsection{Mise à jour, désinstallation, données}\label{sec:donnees}
Les paquets vérifient au démarrage (désactivable) la dernière version publiée sur GitHub et proposent de
l'installer : seule l'archive du code (\texttt{coupole-app-x.y.z.zip}, environ 1 Mo) est téléchargée, extraite par
l'application, l'ancienne version étant restaurée en cas d'échec. Une installation par pip/pipx affiche la commande
à lancer. Le paquet \texttt{.deb} ne se met pas à jour par archive : ses fichiers (\texttt{/opt/coupole}) appartiennent à
dpkg ; Coupole signale la nouvelle version, donne l'adresse du nouveau \texttt{.deb} et la commande \texttt{apt} à
lancer (bouton \emph{Télécharger le paquet}). Les réglages, le cache et les rapports vivent hors du dossier
d'installation et survivent aux mises à jour comme à la désinstallation :

\begin{center}\small
\begin{tabular}{@{}ll@{}}
\toprule
Windows & \texttt{\%LOCALAPPDATA\%\textbackslash Coupole}\\
macOS & \texttt{\textasciitilde/Library/Application Support/Coupole} (cache : \texttt{\textasciitilde/Library/Caches/Coupole})\\
Linux & \texttt{\textasciitilde/.config/coupole} (cache : \texttt{\textasciitilde/.cache/coupole})\\
Tous & variable \texttt{COUPOLE\_HOME} : remplace les deux (installations portables, tests)\\
\bottomrule
\end{tabular}
\end{center}

% =====================================================================
\section{Prise en main}
La fenêtre comporte une barre de modules à gauche (Ctrl+1 à Ctrl+9) et le module courant à droite. Chaque élément
interactif porte une info-bulle. Menu \emph{Aide} : aide de l'écran affiché (F1), ce manuel (Maj+F1), raccourcis
clavier, signaler un problème, rapports d'incident, recherche de mise à jour, À propos (version, licence,
configuration détectée).

\begin{center}\small
\begin{tabular}{@{}ll@{}}
\toprule
F1 & aide de l'écran\\ Maj+F1 & manuel PDF\\ Ctrl+1 … Ctrl+9 & module 1 à 9\\ Ctrl+, & Préférences\\
Ctrl+Maj+A & assistant ASTAP\\ Ctrl+Maj+D & thème sombre / clair\\ Ctrl+R & Banque OHP : rafraîchir l'inventaire\\ Ctrl+Q & quitter\\
\bottomrule
\end{tabular}
\end{center}

\capture{img/reglages_fr.png}{0.62}
Préférences (Ctrl+,) : langue de l'interface ; langue des noms de dossiers, d'objets et des commentaires FITS ;
apparence (section~\ref{sec:apparence}) ; boîtes de dialogue de fichiers (section~\ref{sec:dialogues}) ; PixInsight déjà ouvert (section~\ref{sec:ouvrir}) ; dossier de sortie ; format ; parallélisme ; débit maximal ; mode économe ; recherche de mise à jour ; vérification des nouveautés de la banque au démarrage et sa fréquence maximale (section~\ref{sec:nouveautes}). L'onglet
\emph{Sources} donne les adresses des services (chapitre~\ref{chap:evolutif}). Un fichier de réglages illisible est
mis de côté (suffixe \texttt{.corrompu-…}) et les valeurs par défaut reprennent : l'application démarre toujours.

\subsection{Choisir un fichier ou un dossier, dossier sur un NAS}\label{sec:dialogues}
Les boutons \emph{Parcourir} ouvrent l'explorateur de fichiers du système, avec ses emplacements, favoris et partages
réseau : Explorateur sous Windows, Finder (NSOpenPanel) sous macOS, et sous Linux celui du bureau par le
\textbf{portail XDG} (\texttt{xdg-desktop-portal}) : Dolphin sous KDE Plasma, Fichiers (Nautilus) sous GNOME. Le paquet
autonome embarque sa propre bibliothèque Qt, qui ne peut pas charger le greffon du bureau installé sur le système :
jusqu'à la 0.1.7, il retombait sur le dialogue générique de Qt, en anglais et sans les emplacements du système. Au
démarrage, si l'utilisateur n'a pas déjà posé \texttt{QT\_QPA\_PLATFORMTHEME}, Coupole regarde (sans attendre plus
d'une seconde) si le portail est installé ou en service avec un moteur qui sait choisir des fichiers, et le prend ;
sinon, sous un bureau GTK (GNOME, Cinnamon, Xfce, MATE…), le dialogue GTK ; sinon, le dialogue de Qt, \textbf{en
français} et complété dans sa barre latérale par les dossiers personnels, \texttt{/media}, \texttt{/mnt},
\texttt{/run/media/\$USER}, les partages ouverts dans le gestionnaire de fichiers (\texttt{/run/user/<uid>/gvfs/…}) et
les montages cifs, nfs et sshfs trouvés dans \texttt{/proc/mounts}.

Préférences > \emph{Boîtes de dialogue de fichiers} : \emph{Système} (défaut) ou \emph{Qt}, le second seulement si le
dialogue du système ne s'ouvre pas ou fonctionne mal (portail défaillant). Sous Linux, revenir à \emph{Système} prend
effet au prochain lancement. Tous les textes fournis par Qt (boutons standard, menus contextuels des champs, dialogue
de repli) suivent la langue de l'interface et changent avec elle.

\paragraph{Dossier de sortie sur un NAS.}
\begin{itemize}[nosep]
\item \emph{Linux} : ouvrir d'abord le partage dans Dolphin ou Fichiers (\texttt{smb://serveur/partage}) ; il apparaît
  sous \texttt{/run/user/<uid>/gvfs/smb-share:server=serveur,share=partage} (GNOME ; KDE avec le paquet
  \texttt{kio-fuse}, sans lequel un partage \texttt{smb://} non monté ne peut pas être choisi par une application),
  ou monter le partage en cifs ou nfs (\texttt{/mnt/partage}, conseillé pour de gros volumes ; options par défaut, \texttt{nobrl} inutile : section~\ref{sec:partage}) ;
\item \emph{macOS} : Finder > Aller > Se connecter au serveur (Cmd+K), \texttt{smb://serveur/partage}, puis choisir le
  volume sous \texttt{/Volumes} ;
\item \emph{Windows} : lecteur réseau (\texttt{Z:}) ou chemin UNC \texttt{\textbackslash\textbackslash
  serveur\textbackslash partage}, qu'on peut aussi taper dans le champ du dossier.
\end{itemize}
Ces chemins sont acceptés partout : base d'état lue en lecture seule (URI SQLite d'un chemin UNC écrite avec une
autorité vide), place libre, test d'écriture, chemins longs Windows (\texttt{\textbackslash\textbackslash?\textbackslash
UNC\textbackslash…}), affichage coupé à la ligne ; un partage est reconnu comme tel (UNC, lecteur réseau, cifs, nfs,
sshfs, gvfs, kio-fuse) et le module Qualité y limite ses lectures parallèles.

\subsection{Apparence}\label{sec:apparence}
Coupole a son propre thème, indépendant de celui de l'ordinateur (un système en mode sombre ne mélange plus ses
couleurs à celles de l'application : style Fusion, palette complète, taille de police fixe, contraste WCAG $\geq$ 4,5
pour le texte, les champs, les boutons, les info-bulles et la sélection). Le thème \textbf{sombre} est celui par défaut ;
le thème clair reste au choix. Il se change par le menu \emph{Affichage} > \emph{Apparence} (entrées cochables
\emph{Clair} / \emph{Sombre}, Ctrl+Maj+D pour basculer), ou dans les Préférences : c'est le même réglage, appliqué
immédiatement et enregistré ; les deux restent synchronisés. Un utilisateur qui avait choisi le thème clair le garde :
le réglage enregistré prime sur le défaut. Les captures de ce manuel sont prises dans le thème par défaut.

\subsection{Réglages conservés}\label{sec:conserves}
Coupole retrouve à chaque lancement l'écran tel qu'on l'a laissé. Sont conservés :
\begin{itemize}
\item \textbf{fenêtre} : taille, position, écran, état maximisé ; taille de chaque dialogue (Préférences et son
  onglet, assistant ASTAP, À propos, signalement, aide, correction d'un objet, ajout d'un site) ;
\item \textbf{disposition} : module affiché, onglet de la Banque OHP, position des séparateurs, largeur et ordre des
  colonnes, colonne et sens du tri de chaque tableau ;
\item \textbf{Banque OHP} : recherche, type, télescope, cases \emph{Nouveaux}, \emph{À vérifier}, \emph{À télécharger
  seulement}, \emph{Sans dates douteuses}, nuit et filtre, objets choisis (rétablis dès que l'inventaire est chargé),
  filtres des anomalies et de la carte du ciel ; onglet Traitement : dossier de sortie \emph{dès qu'il est modifié},
  format, langue des noms, \emph{Garder les doublons}, \emph{Garder les FITS}, mode ASTAP, \emph{Vérifier la qualité} ;
\item \textbf{Qualité des images} : dossier analysé, mode échantillon et nombre d'images par lot ;
\item \textbf{Spectres et séries} : dossier du dialogue d'ouverture et menu \emph{Récents} (10 derniers fichiers ;
  un fichier disparu y est grisé) ; axe choisi ;
\item \textbf{Cosmologie} : jeu de paramètres, paramètres personnalisés ($H_0$, $\Omega_m$, gardés même en repassant
  par Planck), $\Omega_k$, $z$ courant, option SH0ES, échelle des courbes (logarithmique ou linéaire) ;
\item \textbf{Sites et heures} : site choisi, zoom et centre de la carte du monde, \emph{Carte en ligne} ;
\item \textbf{dialogues de fichiers} : chacun rouvre dans le dernier dossier utilisé (export CSV, réorganisation,
  ASTAP, signalement…).
\end{itemize}
La disposition vit dans \texttt{interface.json}, à côté de \texttt{reglages.json} (section~\ref{sec:donnees}) ; les
choix de traitement restent dans \texttt{reglages.json}. L'écriture est \emph{différée et groupée} : au plus une toutes
les 2~s et une à la fermeture, atomique (fichier temporaire puis remplacement) ; une frappe dans un champ n'écrit
jamais le disque. Lecture tolérante : une valeur absente, d'un mauvais type ou hors bornes laisse le défaut ; un
fichier illisible est mis de côté (\texttt{.corrompu-…}). Garde-fous : un écran débranché ou une position hors des
écrans recentre la fenêtre, une taille plus grande que l'écran est ramenée à l'écran (comme à 150 ou 200~\%). Un
dossier gardé sur un partage réseau absent ne bloque jamais le démarrage : son existence est vérifiée en fond, et le
dialogue s'ouvre au dossier par défaut tant qu'elle n'est pas confirmée. Rien de ce qui serait dangereux à rejouer
n'est gardé : un traitement interrompu n'est jamais relancé seul (il reprend quand on le relance).

Pour revenir à l'origine : Préférences > \emph{Réinitialiser la disposition} (immédiat ; la langue, le thème, le
dossier de sortie et les autres réglages ne changent pas), ou \cmd{coupole --reinitialiser-interface}
(\cmd{--reset-interface}).

% =====================================================================
\section{Module Banque OHP}
\subsection{La banque}
« OHP student observations » réunit les images prises par les étudiants du DU à l'Observatoire de Haute-Provence :
7~989 images FITS au 8 octobre 2026, 166 objets sous 194 noms, 60 nuits de 2015 à 2026, 78,0 Go ; télescope de
1,20 m (T120, 7~419 images, 1024 pixels de 0,77″, champ de 13,1′) et télescope robotique IRIS de 50 cm (570 images,
2048 pixels de 0,68″ ou 4096 pixels de 0,46″, champs de 23′ et 31′). Les images sont déjà calibrées (dark, flat) et
résolues astrométriquement : \textbf{elles ne doivent pas être recalibrées}.

\subsection{Catalogue}
\capture{img/catalogue_fr.png}{1.0}
L'inventaire vient du service TAP public de l'Observatoire (\texttt{http://tap-ufe.obspm.fr/tap}, table
\texttt{ivoa.obscore}, une requête ADQL synchrone), gardé en cache ; un instantané est livré pour travailler hors
ligne. \emph{Rafraîchir l'inventaire} (Ctrl+R) réinterroge le service et signale les images et les noms nouveaux
(« nouveau depuis le … »). Les noms de la base sont regroupés en objets canoniques ; la sélection d'un ou plusieurs
objets affiche leurs images à droite, filtrables par nuit, filtre et dates douteuses ; cette liste est la sélection
à traiter. L'estimation donne le volume à télécharger, le volume attendu en sortie et la place libre. Les heures
sont affichées en UTC et à l'heure du site (survol : également l'heure de l'utilisateur).

\paragraph{Ce qu'on possède déjà.} D'après le dossier de sortie (\texttt{\_traitement/etat.sqlite}, relu en fond à
l'ouverture, après chaque traitement et à chaque changement de dossier, jamais bloquant), la liste des images
distingue par une pastille et une couleur de texte douce (lisible dans les deux thèmes, contraste $\geq 4{,}5$) :
\textbf{possédée} (convertie et rangée, coche verte), \textbf{doublon écarté} (mêmes pixels qu'une image déjà
convertie, rond gris), \textbf{échec} (à retenter, triangle orange) et \textbf{à télécharger} (texte normal,
flèche) ; l'info-bulle de la pastille donne le statut, la date et le chemin local du fichier. La \textbf{liste des
objets} montre le même état, aux mêmes couleurs : chaque ligne commence par une pastille (coche verte : tout possédé,
doublons écartés compris ; demi-disque vert : en partie ; triangle orange : au moins un échec ; flèche : rien de
téléchargé), le nom de l'objet prend la même couleur, et l'info-bulle détaille (« 12 possédées / 12 · 0 à télécharger
· 0 échec(s) · 3 doublon(s) écarté(s) »). La colonne \emph{possédé} (« 120 / 300 », mini-barre) suit directement le
nom (un ordre de colonnes choisi à la main est respecté) ; la trier range les objets par état (échec, rien, en
partie, complet). La légende, sous les deux listes, vaut pour les deux. La ligne de résumé de l'inventaire ajoute la
possession globale (« 7 989 images (144 doublons), 166 objets, 60 nuits, 78,03 Go — possédées : 7 625 (+ 220
doublons écartés), échecs : 0, à télécharger : 0 »). La case \emph{À télécharger seulement} cache les objets
entièrement possédés et, dans la liste, les images déjà traitées ; une liste vidée l'explique au centre (« Tout est
déjà téléchargé dans … »), comme la liste des images tant qu'aucun objet n'est choisi (« Choisissez un ou plusieurs
objets… » ; au premier lancement, le premier objet est choisi d'office). L'estimation ne compte que ce qui manque
(« 180 images manquantes sur 300 : 1,5 Go »). Changer de dossier de sortie (Préférences, champ ou \emph{Parcourir}
de l'onglet Traitement) relit aussitôt sa possession en fond : pastilles, colonne et résumé suivent, et la barre
d'état l'annonce (« Possession recalculée : 7 625 images trouvées dans … »).
Ligne de commande : \cmd{coupole ohp inventaire --manquantes [--dest DOSSIER] [--json]}, et
colonnes \texttt{possedee}, \texttt{statut\_local}, \texttt{fichier\_local} dans le CSV de \cmd{coupole ohp images}.

\subsection{Traitement}
\capture{img/traitement_fr.png}{1.0}
Pour chaque image, hors doublons :
\begin{enumerate}[leftmargin=*,itemsep=1pt]
\item téléchargement : au plus 4 simultanés (3 par défaut), débit total plafonné (8 Mo/s par défaut), User-Agent
  identifiant Coupole, reprise par requêtes HTTP Range, contrôle de la taille (serveur et inventaire), du multiple
  de 2~880 octets et de la carte \texttt{SIMPLE} ;
\item lecture des cartes FITS telles qu'écrites et des pixels ; empreinte SHA-1 des pixels (doublons exacts) ;
\item contrôle de la solution astrométrique (section~\ref{sec:astrometrie}) ;
\item en-tête corrigé (section~\ref{sec:entetes}) ;
\item écriture au format choisi, relecture et comparaison pixel à pixel ;
\item à la fin : tri en lots, \texttt{LOT.txt}, \texttt{INDEX\_LOTS.csv}, \texttt{\_traitement/journal.csv}.
\end{enumerate}
L'état est gardé dans \texttt{\_traitement/etat.sqlite} : \emph{Arrêter} puis \emph{Lancer} (ou relancer la même
commande) reprend sans rien refaire. Un dossier garde la langue de noms et le format de son premier traitement.
\emph{Garder les doublons} traite aussi ce qui serait écarté ; \emph{Garder les FITS d'origine} les conserve dans
\texttt{\_traitement/fits\_origine}.

La chaîne est un \emph{pipeline} : les téléchargements (fils) et les conversions (processus séparés, un par cœur
disponible) se recouvrent, la file des FITS en attente est bornée ($2\times$ conversions + téléchargements ; si la
place libre à destination est juste, elle est réduite automatiquement — jamais sous conversions + 1 — et le journal le
dit, plutôt que de refuser le traitement), et le
débit global est celui du réseau ou du plafond choisi (mesuré : 24 images de 8 Mo à 8 Mo/s en 24,5 s, soit la durée
du seul téléchargement ; sans plafond, 86 Mo/s de FITS convertis sur 6 cœurs). \emph{Pause} suspend
l'alimentation du pipeline (les images en cours se terminent) et \emph{Reprendre} la rétablit ; \emph{Arrêter}
interrompt téléchargements et conversions sans attendre (les processus sont terminés, ASTAP compris : par signal
sous Linux et macOS, par un \emph{job object} sous Windows). À la fin, un
rapport résume images, durée, converties, doublons, échecs (avec leur raison), volume écrit et lots.

Sur un partage réseau (NAS par SMB), où chaque accès est un aller-retour, l'état est validé au plus une fois par
seconde et \texttt{JOURNAL.txt} écrit par paquets (une coupure ne fait refaire que la dernière seconde) ; le rangement
final ne réécrit que les \texttt{LOT.txt} qui ont changé et ne parcourt plus toute l'arborescence : relancer un
traitement pour une image ne coûte plus une minute de rangement.
La base d'état, elle, n'est jamais écrite directement sur un partage : section~\ref{sec:partage}.

\subsection{Tout télécharger, journal, réorganiser}\label{sec:tout}
\emph{Tout télécharger…} prend la banque entière (environ 8~000 images, 78 Go de FITS, 30 Go en XISF). Au premier
usage (ou si aucun dossier de sortie n'est défini), Coupole demande le dossier de destination par un dialogue natif et
propose \texttt{Documents/Coupole/OHP\_DU\_ECU} ; puis il affiche, \textbf{avant toute confirmation}, le volume à
télécharger, le volume attendu en sortie, la durée estimée au débit plafond (78 Go à 8 Mo/s : 2~h~42) et la place
libre ; il refuse clairement s'il en manque. Le traitement est celui de toute sélection : reprise après coupure ou
fermeture (état par image dans \texttt{\_traitement/etat.sqlite}, FITS déjà téléchargés retrouvés), pause, rapport
de fin. Ligne de commande : \cmd{coupole ohp tout --oui}.

\emph{Ouvrir le journal} affiche \texttt{\_traitement/JOURNAL.txt} : une ligne horodatée (UTC) et bilingue par
événement --- début et fin de session (reprise, pause), chaque image téléchargée, convertie ou écartée avec sa raison,
chaque erreur, chaque conflit de nom ; rotation à 5 Mo (\texttt{JOURNAL-1.txt}, \texttt{JOURNAL-2.txt}). Le détail
tabulaire reste dans \texttt{journal.csv}.

\emph{Réorganiser des fichiers…} range dans l'arborescence des lots des images déjà converties par Coupole qui se
trouvent ailleurs ou selon un ancien rangement : chaque fichier est reconnu par son en-tête (adresse d'origine
\texttt{OHP:Source:URL}, ou carte \texttt{HISTORY} « converti de … »), rattaché à l'inventaire et \textbf{déplacé}
(jamais copié) ; un nom déjà pris par un fichier étranger reçoit un suffixe et une ligne de journal ; rien n'est
jamais écrasé. Les en-têtes sont lus par plusieurs processus, la progression s'affiche (« 1 200 / 7 625 en-têtes
lus ») et \emph{Arrêter} interrompt la lecture (ce qui est déjà reconnu est rangé). Ligne de commande :
\cmd{coupole ohp reorganiser DOSSIER --dest DEST}.

\subsection{Dossier de sortie sur un partage réseau}\label{sec:partage}
\textbf{Pourquoi.} La base d'état (\texttt{\_traitement/etat.sqlite}) est une base SQLite, qui pose des verrous sur
son fichier à chaque écriture. Sous Linux, un partage SMB monté par \texttt{cifs} avec les options par défaut
transmet ces verrous au serveur comme verrous de plage SMB, et l'écriture échoue : « database is locked » au bout
d'une minute, base de 0 octet (vérifié sur un vrai serveur Samba ; c'était le cas de toute version jusqu'à 0.1.9 :
télécharger les nouveautés, \emph{Tout télécharger} ou réorganiser vers un dossier du NAS échouait). La lecture
seule et la copie de fichiers, elles, fonctionnent.

\textbf{Ce que fait Coupole (0.1.10).} Quand le dossier de sortie est reconnu comme un partage (cifs, smb3, nfs,
sshfs, gvfs, kio-fuse sous Linux ; smbfs, afpfs, nfs sous macOS ; chemin UNC ou lecteur réseau sous Windows) ou
quand SQLite n'arrive pas à y écrire en 3~s (partage non reconnu, autre programme), il travaille sur une
\textbf{base de travail locale}, dans le dossier de cache de l'utilisateur (\texttt{bases/<clé>/etat.sqlite},
une par dossier de sortie, la clé étant tirée du chemin), et le dit dans le journal et la barre d'état :
« Dossier sur un partage réseau : base de travail locale, recopiée sur le partage toutes les 30~s ».
\begin{itemize}[nosep]
\item \emph{Au début} : la base du partage est copiée (simple copie de fichier) dans la base de travail.
\item \emph{Pendant le traitement}, toutes les 30~s, \emph{à la fin}, à l'\emph{arrêt} et à l'\emph{annulation} :
  la base de travail est recopiée sur le partage --- instantané cohérent, \texttt{PRAGMA integrity\_check},
  copie dans un fichier temporaire du même dossier, relecture et comparaison de l'empreinte SHA-256, puis
  remplacement d'un coup (\texttt{os.replace}). Aucun verrou n'est jamais posé sur le partage : pas de verrou
  orphelin possible après un plantage, et un lecteur (autre ordinateur, NAS) voit toujours une base complète.
\item \emph{Au démarrage suivant} : si la base du partage est plus récente (traitée ailleurs), elle est recopiée ; si
  la base de travail a des écritures qui n'ont pas été recopiées (plantage, coupure), elles le sont aussitôt.
\item \emph{Deux écrivains} (deux ordinateurs, ou le NAS, ont traité vers le même dossier pendant ce temps) : rien
  n'est écrasé. Coupole le dit clairement et propose une \textbf{fusion} : pour chaque image, le statut le plus
  avancé est gardé (convertie > doublon écarté > échec > en cours), à égalité le plus récent ; empreintes de
  pixels réunies. Les deux bases d'origine sont gardées à côté de la base de travail
  (\texttt{etat.sqlite.local-<date>}, \texttt{etat.sqlite.partage-<date>}). Ligne de commande :
  \cmd{coupole ohp fusionner --dest DOSSIER}.
\end{itemize}
La possession, les nouveautés et les anomalies lisent la base de travail quand elle est à jour (plus rapide qu'une
lecture page par page sur le réseau), sinon la base du partage, en lecture seule. \texttt{JOURNAL.txt}, les CSV et
les \texttt{LOT.txt} restent écrits directement sur le partage (fichiers texte, remplacés d'un coup). Le cache du
module Qualité n'est qu'un cache : il reste dans le dossier de cache de l'utilisateur, sans recopie
(section~\ref{sec:qualite}).

\textbf{Selon le système.} \emph{Linux} : cifs par défaut échoue (cause ci-dessus) ; monter avec \texttt{nobrl}
n'est plus nécessaire. \emph{macOS} : le client smbfs gère ses verrous autrement (selon la façon dont SQLite a été
compilé) ; ce n'a pas pu être vérifié ici, Coupole n'en dépend donc pas et utilise la même base de travail.
\emph{Windows} : SQLite sait en principe écrire sur un partage (verrous SMB natifs), mais plusieurs écrivains y
risquent d'abîmer la base ; même base de travail, vérifiée en intégration continue sur
\texttt{\textbackslash\textbackslash localhost\textbackslash C\$}. Si la base du partage est ouverte par un autre
programme au moment du remplacement, Coupole réessaie puis attend le passage suivant.

\subsection{Nouveautés de la banque}\label{sec:nouveautes}
À l'ouverture (au plus une fois par jour, réglable, désactivable dans les Préférences), si Internet est disponible et
qu'une copie locale existe dans le dossier de sortie, Coupole compare l'inventaire TAP frais à la copie. S'il y a des
images nouvelles, un bandeau apparaît dans la Banque OHP (et un message dans la barre d'état) : « N nouvelles images
(M objets, X Go) depuis le JJ/MM --- Télécharger maintenant ? / Plus tard / Voir ». \emph{Télécharger} traite les
seules nouveautés dans la même arborescence ; \emph{Voir} filtre le catalogue ; rien n'est jamais téléchargé sans
accord explicite. Le bouton \emph{Vérifier les nouveautés} lance la comparaison à la demande. Ligne de commande :
\cmd{coupole ohp nouveautes} (liste) et \cmd{coupole ohp nouveautes --telecharger}.

\subsection{Solution astrométrique}\label{sec:astrometrie}
Contrôles de cohérence, toujours faits : solution présente et lisible ; échelle à moins de 2~\% de la nominale
(\texttt{XPIXSZ}/\texttt{FOCALLEN}) ou de la médiane de la nuit ; angle à moins de 5° de la médiane de la nuit
(modulo 180°) ; centre à moins de 1,5° du médoïde de l'objet (fixe) ou 2° du médoïde objet et nuit (mobile).
Médianes et médoïdes viennent de l'inventaire entier.

Avec ASTAP (facultatif), chaque image est en outre résolue indépendamment :
\begin{center}\small
\begin{tabularx}{\linewidth}{@{}lX@{}}
\toprule
\textbf{Statut} & \textbf{Signification}\\
\midrule
confirmée & ASTAP retrouve le même centre à moins de 10″ et la même échelle à moins de 1~\%.\\
validée & solution gardée : contrôles bons, ASTAP muet (champ pauvre, trop dense, pose courte) ou absent.\\
refaite & solution absente ou incohérente, remplacée par celle d'ASTAP (avec SIP) ; l'ancienne passe en HISTORY.\\
douteuse & contrôles en échec et pas d'ASTAP pour trancher : image rangée à part.\\
échec & rien d'utilisable : image rangée dans \texttt{\_sans\_solution\_astrometrique}.\\
\bottomrule
\end{tabularx}
\end{center}

\subsection{En-têtes corrigés}\label{sec:entetes}
Toutes les cartes d'origine sont gardées dans l'ordre ; chaque valeur remplacée est recopiée dans une carte
\texttt{HISTORY valeur d'origine : …}.
\begin{center}\small
\begin{tabularx}{\linewidth}{@{}>{\ttfamily}lX@{}}
\toprule
\normalfont\textbf{Mot-clé} & \textbf{Correction}\\
\midrule
OBJECT & nom canonique (alias de la base fusionnés).\\
FILTER, FILTSYS, FILTORIG & filtre normalisé (B, V, R, Ha, OIII, g, r, i, z), système (Johnson, Cousins, SDSS, Gunn, Narrowband), valeur d'origine.\\
LATITUDE, LONGITUD & corrigées \textbf{seulement si l'inversion est démontrée} (valeurs échangées retombant sur le site) ; justes : intactes ; ambiguës : signalées, intactes. Opération idempotente.\\
SITELAT, SITELONG, SITEELEV & ajoutés s'ils manquent (site MPC 511 converti en WGS 84) ; présents et cohérents : gardés ; incohérents : signalés.\\
RA, DEC, OBJCTRA, OBJCTDEC & centre réel de la solution.\\
PIXSCALE, ORIENTAT & échelle mesurée, angle de position de l'axe Y.\\
FOCALLEN, XPIXSZ & recalculés seulement s'ils manquent ou sont incohérents de plus de 2~\%.\\
CTYPE1/2 & « -SIP » sans coefficients (IRIS, PinPoint) ramené à TAN.\\
DATE-OBS, DATEDOUT & ajoutée si absente ; \texttt{DATEDOUT = T} quand la date est partagée avec une autre pose.\\
XISFCONV & description de la conversion.\\
\bottomrule
\end{tabularx}
\end{center}
La colonne \texttt{t\_min} de la base vaut \texttt{DATE-OBS} − pose/2 ; \texttt{DATE-OBS} (début de pose) fait foi.
Une date d'en-tête écrite en heure locale (écart entier d'heures égal au décalage du fuseau du site avec
\texttt{MJD-OBS}/\texttt{JD}, ou pose « nocturne » datée en plein jour) est signalée dans le journal, jamais
corrigée en silence.

\subsection{Lots}
\capture{img/lots_fr.png}{0.9}
Un lot réunit ce qui s'empile ensemble : même objet, même instrument (T120, T120-bin1, T120-sCMOS, IRIS-2048,
IRIS-4096), même filtre, et même champ pour un objet fixe (toutes nuits) ou même nuit pour un objet mobile.
\begin{ccode}
07\_Nebuleuses/NGC\_6888/champ\_1\_T120/V/20250716-214427\_NGC\_6888\_V\_1s.xisf\\
01\_Asteroides/914\_Palisana/nuit\_2023-08-15\_T120/R/LOT.txt
\end{ccode}
\texttt{LOT.txt} (français puis anglais) donne objet, noms de la base, instrument, filtre, poses et pose totale,
nuits, centre moyen, angle, statut des solutions, dates douteuses et conseil d'empilement (StarAlignment ou
CometAlignment, rejet adapté au nombre de poses). Noms de fichiers : date et heure UTC du début de pose, objet,
filtre, pose ; tri par nom = tri dans le temps. Noms sûrs sous Windows. L'onglet \emph{Lots} ajoute une colonne
\emph{complet} : poses converties contre poses de la banque pour la même clé (objet, instrument, filtre, et nuit
pour un objet mobile) ; un objet fixe découpé en plusieurs champs compte ses lots ensemble.

\paragraph{Possédé $\neq$ empilable ensemble.} « 840 / 840 » dans la colonne \emph{possédé} dit que toutes les
images d'un objet sont là, pas qu'elles forment un seul empilement : la colonne \emph{lots} de la liste des objets
et son info-bulle donnent la répartition (« 840 images rangées en 9 lots : champ 1 T120 : B 200 · V 300 · R 300 ;
… — on empile lot par lot »), et le clic droit sur un objet propose \emph{Voir les lots de cet objet} (onglet Lots
filtré, bouton \emph{Tous les lots}) et \emph{Ouvrir le dossier de la cible} (double-clic : de même).

\paragraph{Pour PixInsight / N.I.N.A.} Sous le tableau des lots, un encadré donne, pour le lot choisi, ce qu'il
faut saisir dans PixInsight (ImageSolver, panneau \emph{Astrometric solution} de WBPP) ou N.I.N.A. : instrument,
\textbf{focale} (celle qui donne l'échelle mesurée : $f = 206{,}265 \times \mathrm{XPIXSZ} / \mathrm{PIXSCALE}$,
par exemple 7~234,1~mm pour le T120 au lieu des 7~200 de l'en-tête d'origine), \textbf{pixel effectif} (binning
compris : T120 iKon-L en binning 2 : 27~\textmu m), binning, \textbf{échelle} (\si{\arcsecond}\,px$^{-1}$), champ et
centre, chacun avec \emph{Copier}, plus \emph{Tout copier} (aussi au clic droit d'un lot : \emph{Copier les
paramètres d'astrométrie}). Les valeurs sont lues dans l'en-tête des images du lot, jamais codées en dur ; les
colonnes \emph{focale}, \emph{pixel} et \emph{échelle} du tableau se masquent au clic droit sur l'en-tête.
\texttt{LOT.txt} donne les mêmes valeurs. Si un logiciel affiche le double de l'échelle, il a appliqué le binning
deux fois : saisir le pixel non binné (13,5~\textmu m) avec le binning 2, ou le pixel effectif avec le binning 1.

Vérifié dans le code de PixInsight (PJSR : \texttt{AstrometricMetadata.js}, \texttt{BPP-Solver.js}) : ImageSolver
lit les propriétés \texttt{Instrument:Telescope:FocalLength} (m) et \texttt{Instrument:Sensor:XPixelSize}
(\textmu m), sinon \texttt{FOCALLEN} et \texttt{XPIXSZ}, sans multiplier par \texttt{XBINNING} ; WBPP résout le
\emph{master} intégré avec ses métadonnées si elles survivent à l'intégration, sinon avec les valeurs de son panneau,
qui restent celles du dernier instrument utilisé : y saisir focale et pixel (un fichier de réglages WBPP par
instrument évite de les ressaisir). Coupole écrit donc \texttt{FOCALLEN} accordé à l'échelle mesurée (ancienne valeur
en \texttt{HISTORY}) et les propriétés \texttt{Instrument:Sensor:XPixelSize}/\texttt{YPixelSize},
\texttt{Instrument:Camera:XBinning}/\texttt{YBinning}, \texttt{Instrument:Camera:Name},
\texttt{Instrument:Telescope:FocalLength} et \texttt{Aperture} (T120 : 1,2~m). Pour les fichiers déjà convertis :
\cmd{coupole ohp metadonnees DOSSIER --reecrire} (sans \texttt{--reecrire} : simulation) complète les en-têtes
\emph{sans toucher aux pixels} (bloc recopié octet pour octet, relu et comparé ; écriture atomique ; journal
\texttt{\_traitement/metadonnees.csv}) et note focale et pixel dans la base d'état ; \cmd{coupole ohp ranger} récrit
ensuite les \texttt{LOT.txt}. Pour vérifier un master : l'ouvrir, \emph{FITSHeader} ou \emph{Image > Properties},
chercher \texttt{XPIXSZ}, \texttt{FOCALLEN} et \texttt{Instrument:*}.

\subsection{Ouvrir une image}\label{sec:ouvrir}
Double-clic sur une image possédée (liste des images, tableau du module Qualité) : l'application associée du
système. Clic droit : \emph{Ouvrir}, \emph{Ouvrir avec} et \emph{Ouvrir l'emplacement du fichier} (fichier
sélectionné : Explorateur \texttt{/select}, Finder \texttt{open -R}, gestionnaire freedesktop par D-Bus
\texttt{FileManager1.ShowItems} — Dolphin, Nautilus… —, sinon le dossier). \emph{Ouvrir avec} ne propose que les
logiciels d'astronomie trouvés sur la machine (PixInsight, Siril, ASTAP, Aladin, N.I.N.A. : PATH, emplacements
usuels, flatpak, \texttt{/Applications}, Program Files) \emph{et} capables de lire ce fichier d'après la table
vérifiée de la section~\ref{sec:compat} ; les autres apparaissent grisés avec la raison (« ASTAP ne lit que le XISF
non compressé : choisir le format FITS dans le Traitement »). Une image pas encore téléchargée : la barre d'état le
dit. Une entrée désactivée porte son motif dans son libellé (« Ouvrir avec (rien de téléchargé) ») et le motif
complet s'affiche dans la barre d'état au survol : plus d'info-bulle par-dessus le menu (sous KDE, elle en captait
la souris).

\paragraph{Comment chaque logiciel est lancé (0.1.12).} Toujours le programme détecté, avec une \emph{liste}
d'arguments (jamais d'interpréteur de commandes) ; chaque lancement est noté dans \texttt{coupole.log} (méthode,
programme, arguments, numéro de processus ; si le programme s'arrête dans les 3~premières secondes : code de retour
et début de ce qu'il a écrit) et son résultat s'affiche dans la barre d'état : « Ouverture dans Siril… », ou
l'erreur (programme introuvable, arrêt immédiat avec son message). Jamais de silence.
\begin{itemize}
  \item \textbf{PixInsight} — Linux : le script \texttt{bin/PixInsight.sh} (l'éditeur précise que le binaire seul
  ne démarre pas sous Linux). Ce script reconstruit ses arguments par un \texttt{eval} qui n'entoure de guillemets
  que ceux qui contiennent une espace : un chemin avec \texttt{\$}, \texttt{`}, \texttt{'}, \texttt{"},
  \texttt{\textbackslash}, \texttt{(}, \texttt{*}… serait interprété ; Coupole lui passe alors un \emph{lien
  symbolique temporaire} au nom sans caractère spécial (pour enregistrer : « Enregistrer sous » dans le dossier
  d'origine). Windows : \texttt{bin\textbackslash PixInsight.exe}. macOS : \texttt{open -a PixInsight.app
  \emph{fichier}}, ou \texttt{open -n -a PixInsight.app --args -n \emph{fichier}} pour une nouvelle instance.
  \item \textbf{PixInsight déjà ouvert.} Sans option, PixInsight confie le fichier à l'instance ouverte (« Yielded
  execution to running application instance \#1 », code~0) ; si celle-ci est occupée (script en cours), \emph{rien
  ne s'ouvre}. D'où le réglage \emph{Préférences > PixInsight déjà ouvert} : « Ouvrir une nouvelle fenêtre (-n) »
  (par défaut, sûr) ou « Envoyer à la fenêtre ouverte (essai) ». Quand une instance tourne (processus détecté), le
  menu propose les deux : « PixInsight (nouvelle fenêtre) » et « PixInsight (fenêtre ouverte, essai) », celle du
  réglage d'abord ; si la sortie dit que PixInsight a cédé la main, la barre d'état le signale.
  \item \textbf{Siril} : \texttt{siril \emph{fichier}} ; en flatpak : \texttt{flatpak run --file-forwarding
  org.free\_astro.siril @@ \emph{fichier} @@} (le portail de documents ne transmet le fichier que si le bac à sable ne
  le voit pas). Le paquet Flathub voit tout le système (\texttt{--filesystem=host}) sauf \texttt{/tmp}, \texttt{/run}
  (hors \texttt{/run/media}), \texttt{/usr}… ; si un réglage \texttt{flatpak override} lui retire un dossier, la
  barre d'état le dit et donne la commande : \cmd{flatpak override --user --filesystem=/mnt/partage
  org.free\_astro.siril}.
  \item \textbf{ASTAP, Aladin} : \texttt{\emph{programme} \emph{fichier}}.
  \item \textbf{N.I.N.A.} n'ouvre pas une image passée en argument (ses seules options : profil, séquence, débogage,
  accélération graphique) : il reste grisé dans \emph{Ouvrir avec}, avec la raison ; ouvrir l'image depuis N.I.N.A.
\end{itemize}

\paragraph{Double-clic qui ouvre PixInsight sans l'image (Linux).} \texttt{xdg-open} lance l'application associée
au type du fichier par son fichier \texttt{.desktop}. Si sa ligne \texttt{Exec=} n'a pas de code \texttt{\%F}
(\texttt{\%f}, \texttt{\%u}, \texttt{\%U}), l'application démarre \emph{sans} le fichier — cas d'un
\texttt{pixinsight.desktop} écrit à la main. Coupole le détecte (\texttt{xdg-mime query filetype}, puis
\texttt{xdg-mime query default}, lecture du \texttt{.desktop}) et lance alors le logiciel directement (barre
d'état : « PixInsight lancé directement… ») ; une autre application dans ce cas n'est pas lancée, un message dit
pourquoi. Pour corriger le bureau tout entier (Coupole ne modifie pas ce fichier), dans
\texttt{\textasciitilde/.local/share/applications/pixinsight.desktop} :
\begin{verbatim}
Exec=/opt/PixInsight/bin/PixInsight.sh %F
MimeType=image/x-xisf;image/fits;
\end{verbatim}
puis \cmd{update-desktop-database \textasciitilde/.local/share/applications} et, au besoin,
\cmd{xdg-mime default pixinsight.desktop image/x-xisf}.

\paragraph{Copie d'un ancien traitement.} Une sortie produite avant Coupole (script \texttt{ohp\_xisf.py}) ou sur
une autre machine note dans sa base d'état le chemin absolu d'alors (« /srv/ancien/OHP\_DU\_ECU/… »).
Coupole retrouve chaque fichier dans le dossier de sortie courant, dans cet ordre : chemin relatif noté dans la base,
\texttt{\_traitement/journal.csv} (colonne \texttt{destination}), ancienne racine ou dossier de type, puis le nom
attendu dans le dossier du lot ; les chemins trouvés par le journal sont notés une fois pour toutes dans la base
(en fond, au premier chargement ; par la base de travail locale sur un partage). Ouvrir, \emph{Ouvrir avec},
l'emplacement, le dossier de la cible et les lots fonctionnent donc sur la copie de la banque entière (7~625 images
retrouvées en 0,34~s, notées dans la base en 0,8~s). \cmd{coupole ohp metadonnees DOSSIER --reecrire} note aussi
ces emplacements.

\subsection{Anomalies}
\capture{img/anomalies_fr.png}{1.0}
\begin{center}\small
\begin{tabularx}{\linewidth}{@{}>{\ttfamily}llX@{}}
\toprule
\normalfont\textbf{Genre} & \textbf{Action} & \textbf{Signification}\\
\midrule
meme\_fichier & écartée & même pose sous un autre chemin ou une autre orthographe\\
copie\_diurne & écartée & copie datée en plein jour d'une image bien datée\\
pixels\_identiques & écartée & mêmes pixels (SHA-1) qu'une autre image, au traitement\\
memes\_meta & signalée & mêmes date, filtre, pose, taille, nom différent\\
date\_partagee & signalée & DATE-OBS identique à celle d'une autre pose\\
date\_diurne & signalée & Soleil au-dessus de l'horizon au site (jamais pour le Soleil ni la radio)\\
champ\_incoherent & signalée & centre loin des autres poses de l'objet\\
classement & signalée & nom classé automatiquement ou non classé\\
instrument & signalée & instrument absent de la table\\
\bottomrule
\end{tabularx}
\end{center}
Au 8 octobre 2026 : 133 + 11 = 144 doublons écartés, 1~654 dates partagées, 3 poses diurnes, 92 champs
incohérents. Export CSV aux colonnes bilingues ; rien n'est supprimé.

\subsection{Carte du ciel}
\capture{img/ciel_fr.png}{0.95}
Projection d'Aitoff, ascension droite croissante vers la gauche, 12~h au centre ; un point par objet fixe (taille :
nombre d'images) et par nuit pour les objets mobiles ; écliptique en tirets, plan galactique en pointillés.

\subsection{Fiche en ligne (facultative)}\label{sec:fiche}
\capture{img/fiche_fr.png}{0.95}
Si Internet est disponible, l'onglet \emph{Fiche en ligne} (ou le bouton du même nom sous le catalogue) affiche des
informations publiques sur l'objet choisi :
\begin{itemize}
\item ciel profond et étoiles : \textbf{SIMBAD} (CDS, Strasbourg) --- type, coordonnées ICRS, magnitudes, parallaxe
(et la distance qu'elle donne quand elle est assez précise), distance mesurée préférée avec sa méthode et sa
référence, vitesse radiale et redshift, taille angulaire, types spectral et morphologique, autres noms ; liens
vers la page SIMBAD, \textbf{Aladin Lite} (image du ciel centrée sur l'objet) et, pour les galaxies, \textbf{NED} ;
si SIMBAD ne connaît pas le nom, \textbf{Sesame} (CDS) interroge aussi NED et VizieR ;
\item astéroïdes, comètes, transneptuniens : base des petits corps du \textbf{JPL} (SBDB) --- classe orbitale,
géocroiseur ou non, éléments (a, e, i, q, Q, $\Omega$, $\omega$, période, époque), diamètre, albédo, période de
rotation, magnitude absolue ;
\item planètes et désignations NEOCP : pas de fiche (message explicatif).
\end{itemize}
\textbf{L'identifiant retenu est toujours affiché} : SIMBAD rattache par exemple « NGC 6888 » à son étoile
Wolf-Rayet HD 192163 ; la fiche le signale. Un redshift positif s'envoie au module Cosmologie (bouton \emph{Envoyer
ce redshift au module Cosmologie}).

\paragraph{Politesse et hors ligne.} Rien ne part tant que l'onglet n'est pas affiché ; une seule requête par
objet choisi (après 0,6~s sans autre changement de sélection), au plus une par seconde et par service, délai de
8~s, User-Agent de Coupole. Les réponses sont gardées 30~jours dans le cache de Coupole (\texttt{enligne.json},
avec leur date ; un « introuvable » un jour). Hors ligne, la dernière fiche connue s'affiche avec sa date, sinon un
message clair ; le reste de l'application n'en dépend jamais. La case \emph{Interroger les services en ligne} coupe
tout envoi. Les adresses (\texttt{simbad.base}, \texttt{simbad.page}, \texttt{aladin.page}, \texttt{ned.page},
\texttt{sbdb.page}, \texttt{resolution.sbdb}, \texttt{resolution.sesame}) se règlent comme les autres sources.

% =====================================================================
\section{ASTAP (facultatif)}
Sans ASTAP, Coupole fait tout le reste. ASTAP (Han Kleijn, licence MPL~2.0, \url{https://www.hnsky.org/astap.htm})
ajoute la vérification indépendante de chaque solution. Il faut le programme et un catalogue d'étoiles.

\capture{img/astap_fr.png}{0.72}

\paragraph{Catalogue.} Le site officiel indique que pour un champ de 0,6° ou plus, D05, D20, D50 et D80 conviennent,
et que les catalogues denses servent les petits champs. Les champs de la banque vont de 10,7′ (0,18°) à 31,4′
(0,52°) : \textbf{D80} est conseillé (environ 1,2 Go ; catalogue du traitement de référence, 6~888 solutions
confirmées sur 7~625) ; D50 (environ 0,9 Go) convient aussi pour IRIS.

\begin{center}\small
\begin{tabularx}{\linewidth}{@{}lX@{}}
\toprule
\textbf{Système} & \textbf{Programme (liens de la page officielle)}\\
\midrule
Windows x64 & \texttt{astap\_setup.exe} puis \texttt{d80\_star\_database.exe} (même dossier, \texttt{C:\textbackslash Program Files\textbackslash astap}) ; facultatif : \texttt{astap\_command-line\_version\_win64.zip}\\
Windows ARM64 & \texttt{astap\_command-line\_version\_win11\_aarch64.zip} dans \texttt{C:\textbackslash astap} ; catalogue : \texttt{d80\_star\_database.exe} (installeur x86, émulé) vers ce dossier, ou \texttt{d50\_star\_database.zip} (pas de D80 en zip)\\
macOS Apple Silicon & \texttt{astap\_M1.pkg}, \texttt{d80\_star\_database.pkg} ; une fois : \texttt{codesign --force -s - /Applications/ASTAP.app/Contents/MacOS/astap}\\
macOS Intel & \texttt{astap.pkg}, \texttt{d80\_star\_database.pkg}\\
Linux Debian/Ubuntu & \texttt{astap\_amd64.deb} (ou \texttt{astap\_aarch64.deb}), \texttt{d80\_star\_database.deb} (dans \texttt{/opt/astap})\\
Linux Fedora/openSUSE & \texttt{astap\_amd64.rpm} ; D80 extrait du \texttt{.deb} (ci-dessous) ou D50 en zip dans \texttt{/opt/astap}\\
Linux Arch/Manjaro & \texttt{astap\_amd64\_gtk3.pkg.tar.zst} (ou GTK2) ; D80 : paquet AUR \texttt{d80-star-db-astap} ou extrait du \texttt{.deb} ; D50 : \texttt{d50\_star\_database.pkg.tar.zst}\\
Autres Linux & \texttt{astap\_amd64.tar.gz} ou \texttt{astap\_command-line\_version\_Linux\_amd64.zip} ; pas d'AppImage officielle\\
\bottomrule
\end{tabularx}
\end{center}

\paragraph{Pas de D80 en zip.} Le lien « D80 (zipped) » de la page officielle est mort (404, vérifié le 9/10/2026) :
le D80 n'existe qu'en \texttt{.exe}, \texttt{.pkg} et \texttt{.deb}. Hors Debian, l'extraire du \texttt{.deb} sans
toucher aux droits de \texttt{/} (commande vérifiée sur le vrai paquet : 1~213~392~756 octets, 1~476 tuiles) :
\begin{ccode}
bsdtar -xf d80\_star\_database.deb data.tar.xz\\
sudo mkdir -p /opt/astap\\
sudo tar -xJf data.tar.xz -C /opt/astap --strip-components=3 --no-same-owner\\
\hspace*{2em}--wildcards \textquotesingle./opt/astap/d80\_*\textquotesingle
\end{ccode}
(jamais \texttt{sudo tar -x … -C /} : l'archive contient \texttt{./} et \texttt{./opt/}). \textbf{Arch/Manjaro} : le
paquet officiel met \texttt{astap} et \texttt{astap\_cli} dans \texttt{/opt/astap} seulement (rien dans le PATH :
Coupole les y trouve) ; l'\texttt{astap} graphique GTK2 demande \texttt{libgtk-x11-2.0}, absent des dépôts officiels
(AUR seulement) et non déclaré ; \texttt{astap\_cli} est autonome (lié statiquement) et suffit à Coupole ; pour
l'interface d'ASTAP, le paquet GTK3. Un test de l'intégration continue vérifie chaque lien, pour chaque système.

\begin{sloppypar}\noindent\textbf{Détection.} Ordre : chemin choisi (Préférences, \cmd{coupole astap --definir}), variable
\texttt{COUPOLE\_ASTAP}, \texttt{astap\_cli} partout (PATH puis emplacements par défaut), puis seulement \texttt{astap}, emplacements par défaut (Windows
\texttt{Program Files\textbackslash astap}, \texttt{\%LOCALAPPDATA\%\textbackslash astap} ; macOS
\texttt{/Applications/ASTAP.app}, \texttt{/usr/local/opt/astap} ; Linux \texttt{/opt/astap}, \texttt{/usr/bin},
\texttt{/usr/local/bin}, \texttt{\textasciitilde/astap}). Catalogue : dossier choisi, \texttt{COUPOLE\_ASTAP\_CATALOGUE},
dossier du programme, \texttt{/opt/astap}, \texttt{/usr/share/astap/data}, \texttt{/usr/local/opt/astap}. Un catalogue
est déclaré complet avec ses 1~476 tuiles ; la version est lue sur \texttt{astap\_cli -h}.\end{sloppypar}

% =====================================================================
\section{Formats de sortie}
\begin{center}\small
\begin{tabularx}{\linewidth}{@{}lXl@{}}
\toprule
\textbf{Format} & \textbf{Contenu} & \textbf{Lecteurs}\\
\midrule
XISF (défaut) & Float32 en ADU, \texttt{bounds="-1000:65535"} ; UInt16 pour les FITS IRIS entiers ; compression zstd+sh niveau 9 ; propriétés \texttt{Observation:*}, \texttt{Instrument:*} & PixInsight\\
XISF compatible & UInt16 + piédestal de 1~000~ADU (\texttt{PEDESTAL}), zlib+sh (ce qu'écrit N.I.N.A.) & PixInsight, Siril $\geq$ 1.4, N.I.N.A.\\
FITS \texttt{.fits.fz} & compression par tuiles sans perte : GZIP\_2 sans quantification (flottants), RICE\_1 (entiers) & Siril, astropy, DS9\\
FITS & float32 (ou UInt16) non compressé & tous\\
\bottomrule
\end{tabularx}
\end{center}
\textbf{XISF} : \texttt{bounds} fixe la plage que PixInsight ramène à [0,1] à l'ouverture : rien à régler, même
échelle pour toutes les poses ; la borne basse à −1~000 garde les valeurs négatives des poses courtes calibrées.
Taille : 36~\% du FITS d'origine pour le T120 (float64), 62~\% pour IRIS. \textbf{FITS compressé} : mesuré avec
astropy 8.0, RICE\_1 et HCOMPRESS\_1 quantifient les flottants (écart d'un ADU, même avec
\texttt{quantize\_level=0}) ; seuls GZIP\_1 et GZIP\_2 sans quantification restituent les flottants bit à bit.
\textbf{FITS float32} : dans PixInsight, régler la plage de lecture FITS sur 0–65~535 avec \emph{Truncate}.
Pour tous les formats sauf le XISF compatible, l'écart au FITS d'origine reste sous le demi-ULP du float32
(vérifié pixel par pixel).

\paragraph{XISF compatible (N.I.N.A., Siril).} UInt16 = ADU + 1~000 arrondis, écrêtés à [0, 65~535] ; le piédestal
garde les fonds négatifs des poses courtes calibrées et figure dans l'en-tête (\texttt{PEDESTAL}, convention de
PixInsight). Mesuré sur 66 poses de la banque (\texttt{outils/audit2/mesure\_perte\_uint16.py}) : arrondi
$\le 0{,}5$~ADU, moyenne du fond changée de 0,014~ADU au plus (bruit du fond : 10 à 109~ADU) ; pixels écrêtés en
haut : 0,012~\% au plus (étoiles au-delà de 64~535~ADU) ; plus aucun pixel perdu en bas. Sans piédestal, une pose
T120 de 10~s au fond sur-soustrait (−220~ADU) perdait 99,75~\% de ses pixels à 0 et 6,1~\% des pixels des poses de
1–2~s tombaient à 0 : d'où le piédestal. Les IRIS entiers restent sans perte.

\subsection*{Qui lit quoi (vérifié)}\label{sec:compat}
\begin{center}\small
\begin{tabularx}{\linewidth}{@{}>{\raggedright}p{0.25\linewidth}ccccX@{}}
\toprule
\textbf{Fichier} & \textbf{PixInsight} & \textbf{Siril} & \textbf{N.I.N.A.} & \textbf{ASTAP} & \textbf{Remarque}\\
\midrule
XISF Float32 (défaut) & oui & non & non & non & Siril $\geq$ 1.4 le lit mais hors de sa plage [0, 1] ; N.I.N.A. : image blanche\\
XISF compatible (UInt16, zlib) & oui & $\geq$ 1.4.0 & oui & non & ASTAP : XISF non compressé seulement\\
XISF UInt16 zstd (IRIS entiers) & oui & $\geq$ 1.4.0 & non (3.2) & non & zstd lu par N.I.N.A. depuis janvier 2026 (développement)\\
FITS Float32 & oui & oui & non & oui & N.I.N.A. suppose [0, 1] ; Aladin, astropy : oui\\
FITS UInt16 & oui & oui & oui & oui & \\
\texttt{.fits.fz} Float32 (GZIP) & — & oui & — & non & ASTAP : Rice seulement\\
\texttt{.fits.fz} UInt16 (Rice) & — & oui & — & $\geq$ 2026.07.16 & — : non vérifié\\
\bottomrule
\end{tabularx}
\end{center}
Sources : \textbf{PixInsight}, code de PCL (\texttt{XISF.cpp} : codecs zlib, lz4, lz4hc, zstd, avec ou sans
\texttt{+sh} ; \texttt{XISFReader.cpp} : \texttt{bounds} obligatoire pour un flottant, écrêtage puis ramené à
[0, 1]) ; spécification XISF 1.0 révision 1 (septembre 2026, \url{https://pixinsight.com/doc/docs/XISF-1.0-spec/XISF-1.0-spec.html}),
schéma XSD identique à celui qu'emploie Coupole. \textbf{Siril}, essai réel des AppImage officielles (siril-cli,
pixels relus et comparés) : 1.2.0 et 1.2.6 ne lisent pas le XISF (« file type not supported »), 1.4.0 et 1.4.4
lisent tous nos codecs bit à bit ; mais les flottants XISF y restent en ADU (statistiques multipliées par 65~535),
alors que les FITS flottants sont ramenés à [0, 1] ; paquets Debian 13 et Ubuntu 24.04 (Siril 1.2.x) : pas de XISF.
\textbf{N.I.N.A.}, code (étiquette \texttt{Version-3.2}, 27/11/2025 : \texttt{XISF.cs} l.~304–328 : lz4, lz4hc,
zlib, avec ou sans \texttt{+sh} ; \texttt{Float32Converter.cs} l.~23 : $(\text{ushort})(v \times 65535)$,
\texttt{bounds} jamais lu) et décodage réel d'un de nos fichiers. \textbf{ASTAP}, page officielle (XISF « uncompressed »,
\texttt{.fits.fz} Rice depuis 2026.07.16, GZIP non). Fichiers d'essai pour les testeurs : \texttt{tests/donnees/}
(image de synthèse en XISF par défaut, XISF compatible et FITS, moins de 300~Ko chacun).

\paragraph{Écart à la révision 1.} Coupole écrit dans \texttt{Observation:Center:RA/Dec} le centre de la solution
astrométrique (choix du traitement de référence, antérieur à la révision) ; la révision 1 réserve ces propriétés au
point visé par le télescope et normalise la solution dans l'espace \texttt{AstrometricSolution:} (non écrit : la
solution reste dans les mots-clés WCS, que PixInsight lit). Gardé tel quel en 0.1.9 pour ne pas changer les fichiers
par rapport au traitement de référence.

% =====================================================================
\section{Qualité des images (facultatif)}\label{sec:qualite}
\capture{img/module_qualite_fr.png}{1.0}
Le module demande la bibliothèque SEP (Source Extractor en Python, LGPL, environ 1 Mo) ; il n'est jamais lancé
automatiquement (bouton \emph{Mesurer}, case \emph{Vérifier la qualité des images} du traitement, ou
\cmd{coupole qualite DOSSIER --ecrire}). Rapport par lot : \texttt{QUALITE.csv} et \texttt{QUALITE.txt}.

\paragraph{Poses de ciel seulement.} Par défaut, seules les poses de ciel (lights) sont mesurées. Sont exclus sans
être ouverts : les dossiers Flat(s), Dark(s), Bias, Offset(s), DarkFlat(s), Calibration, Master(s) (« FLAT L »,
« Dark-Bias » compris), cosmetized, calibrated, registered, debayered, \texttt{\_platesolve…} et leurs noms
français (Plats, Noirs, Biais) ; les fichiers \texttt{FLAT\_}, \texttt{DARK\_}, \texttt{BIAS\_},
\texttt{DARKFLAT\_} (N.I.N.A., ASIAIR), les masters (\texttt{masterDark}, \texttt{masterFlat},
\texttt{masterLight}, \texttt{superbias}) et les sorties de WBPP (\texttt{\_c}, \texttt{\_cc}, \texttt{\_r}…).
En dernier recours, une pose dont l'en-tête (\texttt{IMAGETYP}, \texttt{FRAME} ou la propriété XISF
\texttt{Observation:Image:Type}) dit flat, dark, bias, offset ou master est exclue à la mesure ; Light, Light Frame,
Science ou Object sont mesurées. Le résumé dit « 70 image(s) trouvée(s), 66 pose(s) de ciel mesurée(s), 0 en erreur,
4 fichier(s) de calibration exclu(s) » ; \emph{Fichiers exclus…} en donne la liste avec le motif. Case
\emph{Inclure aussi les poses de calibration} (ligne de commande : \cmd{--avec-calibration}) : tout est mesuré.

\paragraph{Fichiers des autres logiciels, erreurs.} Les XISF de N.I.N.A. (mots-clés sans commentaire), de l'ASIAIR
(sans \texttt{Metadata}), de PixInsight et de Siril sont lus par un lecteur tolérant : seul ce qui empêche
vraiment de lire les pixels est refusé ; nos propres fichiers restent contrôlés strictement. Un fichier non mesurable
apparaît dans le tableau, colonne \emph{état}, avec son motif (« erreur : format XISF non lu (…) ») ; le résumé
compte les erreurs et \texttt{coupole.log} garde la pile ; la ligne de commande sort avec le code 1 si aucune image
n'a pu être mesurée. Pose N.I.N.A. d'une caméra CMOS (9~576 × 6~388, 61~Mpx) : 4,2~s, 1~170~Mo de pointe (calcul en
float32), FWHM 3,5″ à 1,695″/px ; le nombre de processus est réduit d'après la taille des images (moitié de la
mémoire libre) et le nombre d'étoiles ajustées croît avec le capteur (120 à 400).

\paragraph{Vitesse, reprise, échantillon.} Les mesures tournent dans des processus séparés (nombre donné par le plan
machine : cœurs, mémoire, mode économe et bridage des Préférences ; au plus 3 quand le dossier est sur un partage
réseau, détecté et signalé, car la lecture limite alors la vitesse), la lecture d'une image recouvrant le calcul des
autres ; l'interface reste libre, \emph{Arrêter} est immédiat. Chaque image mesurée est gardée aussitôt dans le
cache \texttt{\_traitement/qualite.sqlite} par (chemin, taille, date) et affichée ; le \texttt{QUALITE.csv} de son
lot (écriture atomique) est réécrit au plus toutes les 5~s et à la fin du lot ou de l'arrêt : relancer ne refait pas
ce qui est fait, fermer puis rouvrir reprend.

\paragraph{Inventaire en flux.} La mesure commence dès le premier lot trouvé, pendant que l'inventaire continue
(« Inventaire du dossier : 1 234 fichiers trouvés dans 56 lots… », \emph{Arrêter} immédiat) ; l'inventaire n'ouvre
aucun fichier et lit taille et date dans la liste de chaque dossier (un \texttt{stat} au plus par image). Dans une
sortie de Coupole, seuls les dossiers des lots sont lus, d'après \texttt{INDEX\_LOTS.csv} ou la base d'état ; les
restes de traitement (\texttt{converties}, \texttt{telechargements}) ne sont jamais pris pour des lots. Sur un
partage, 32 fils lisent les dossiers. Le cache n'est \textbf{jamais} écrit sur un partage réseau : le client SMB de
Linux transmet les verrous de SQLite au serveur, qui les refuse (« database is locked » au bout de 30~s, trois fois
avant la première mesure en 0.1.8 : les 90~s d'« Inventaire du dossier… » observées sur un vrai partage) ; il va dans
le dossier de cache de l'utilisateur (un ancien cache du partage est relu). Mesuré sur un vrai partage Samba, 831
XISF de 09\_Galaxies : 0.1.8 92,7~s avant la première mesure (latence de 10~ms) ; 0.1.9 inventaire complet en
0,24–0,57~s (10~ms) et 0,57–1,05~s (20~ms), premier résultat en 0,2–0,6~s. Au-delà de 200 images, Coupole propose un \textbf{échantillon} de 5 images par lot (réparties dans
le temps : première, dernière, milieu…, réglable dans l'onglet et par \cmd{--echantillon N}), qui suffit à juger un
lot (mesuré en premier, pendant que vous choisissez), ou \emph{Tout mesurer} avec la durée estimée d'après les
images déjà mesurées ; la progression donne le temps
restant et le débit (images/min). Sur un processeur à 6 cœurs : 1,21 → 0,33~s par image T120
(modèle de Moffat vectorisé, dérivées analytiques, 120 étoiles par image), soit la banque entière en une dizaine de
minutes sur 5 processus.

\begin{center}\small
\begin{tabularx}{\linewidth}{@{}lXl@{}}
\toprule
\textbf{Mesure} & \textbf{Méthode} & \textbf{Validation synthétique}\\
\midrule
FWHM & Moffat elliptique à $\beta$ libre ajusté sur chaque étoile isolée et non saturée (modèle intégré sur le pixel) ; médiane & ≤ 0,4~\% (Gauss et Moffat, 2,5–8 px)\\
ellipticité & $1-b/a$ du même ajustement & ± 0,02\\
carte 3 × 3 & médianes par neuvième de champ & ordre retrouvé\\
fond, bruit & SEP, maille de 64 px & 0,5~\% ; 5~\%\\
gradient, résidu & plan ajusté sur la carte du fond ; amplitude du résidu & gradient et vignetage retrouvés\\
saturées & pic au-delà de SATURATE/DATAMAX (sinon 60~000 ADU) & 5 sur 5\\
traînées & objets de grand axe ≥ 12 px et $a/b \ge 6$ & 1 sur 1, 0 faux\\
RSN & seulement avec le gain (EGAIN/GAIN) & donné ou absent\\
échantillonnage & FWHM < 2 px, 2–5 px, > 5 px & —\\
\bottomrule
\end{tabularx}
\end{center}
Sur 14 images IRIS (12 nuits, FWHM de 2,4 à 11,3 px), la FWHM mesurée et celle écrite par PinPoint dans l'en-tête
sont corrélées (r = 0,974), rapport médian 1,046 (0,85 à 1,16) : les deux estimateurs diffèrent (méthode de PinPoint
non décrite). Non proposés faute de validation : mag\,arcsec$^{-2}$ (pas de point zéro), détection des « donuts » de
poussière. Une FWHM qui croît vers un bord peut venir du capteur, de la courbure de champ, de la collimation, mais
aussi de la turbulence, du suivi ou de la mise au point : la carte seule ne tranche pas.

% =====================================================================
\section{Spectres et séries}
\capture{img/module_donnees_fr.png}{1.0}
Formats lus : FITS (image 1D avec axe WCS linéaire \texttt{CRVAL/CDELT/CRPIX}, cubes radio 1 × 1 × N, tables
binaires) et CSV ou texte en colonnes (en-tête facultatif, unité entre parenthèses ou crochets). Vitesse radio :
$v = c\,(1 - f/f_0)$ (convention radio de la norme FITS), $f_0$ lue dans \texttt{RESTFRQ}, sinon raie H\,I à
1~420,405~751~768 MHz. Le référentiel est celui de l'axe des fréquences (\texttt{SPECSYS}) ; aucune correction LSR
ou barycentrique n'est faite implicitement. Les formats des radio-antennes d'enseignement s'ajoutent par un lecteur
(CONTRIBUTING.md, « Ajouter un format »).

% =====================================================================
\section{Sites et heures}
\capture{img/module_sites_fr.png}{0.95}
Le temps est toujours tenu en UTC (dates « aware ») ; l'heure locale du site et celle de l'utilisateur sont
affichées en plus. La « date du soir » d'une nuit est la date locale du site à midi précédent ; le plein jour se
calcule par la hauteur du Soleil au site (astropy, sans téléchargement IERS). Sites livrés : OHP (MPC 511), Meudon
(005), Paris (007), coordonnées géodésiques calculées depuis les constantes de parallaxe du Minor Planet Center.
Les sites partenaires s'ajoutent (fuseau IANA, code MPC s'il existe) quand leurs coordonnées sont connues de
source sûre. Carte : tuiles OpenStreetMap en cache disque, deux téléchargements au plus, User-Agent identifiant
l'application, attribution affichée ; hors ligne, fond simple et points. Sur un écran dense (mise à l'échelle
$\geq 150$~\%, Mac Retina), la carte dessine les tuiles du niveau de zoom supérieur à demi-taille (quatre par tuile
logique, l'équivalent « @2x ») : elle reste nette au lieu d'être agrandie par le système.

% =====================================================================
\section{Cosmologie}\label{sec:cosmo}
\capture{img/module_cosmo_fr.png}{1.0}
Le redshift $z$ d'un objet lointain ne donne pas « une » distance mais plusieurs, toutes justes, qui répondent à des
questions différentes. Le module les calcule dans le modèle $\Lambda$CDM avec
\[ E(z)^2 = \Omega_r(z)(1+z)^4 + \Omega_m(1+z)^3 + \Omega_k(1+z)^2 + \Omega_\Lambda, \qquad
   D_C = \frac{c}{H_0}\int_0^z\frac{dz'}{E(z')}, \]
rayonnement compris (photons du fond diffus et neutrinos, dont un massif de 0,06~eV).
\begin{center}
\begin{tabularx}{\linewidth}{@{}lX@{}}
\toprule
Grandeur & Usage \\
\midrule
$D_C$ comobile radiale & distance actuelle de l'objet ; sature à $\approx 46$~G\,al (horizon des particules)\\
$D_M$ comobile transverse & égale à $D_C$ si l'univers est plat ; entre dans $D_L$ et $D_A$\\
$D_L = (1+z)D_M$ & photométrie : $F = L/(4\pi D_L^2)$ ; module de distance $\mu = 5\log_{10}(D_L/10~\mathrm{pc})$\\
$D_A = D_M/(1+z)$ & tailles : $\theta = \ell/D_A$ ; échelle en kpc par seconde d'arc ; maximum vers $z\approx1{,}6$\\
$c\,t_L$ & chemin parcouru par la lumière (le plus intuitif, le moins physique)\\
$t_L$, âge à $z$, $t_0$ & temps de regard en arrière ; contrôle $t_L + t(z) = t_0$\\
$V_C$ & volume comobile jusqu'à $z$ (Hogg 1999, exact aussi en univers courbe)\\
vitesses & $cz$ (naïf), Doppler relativiste (inadapté), $H_0 D_C$ (la seule correcte)\\
\bottomrule
\end{tabularx}
\end{center}

\paragraph{Paramètres.} Planck 2018 par défaut ($H_0 = 67{,}66$, $\Omega_m = 0{,}3111$ dont $\Omega_\nu = 0{,}00144$,
$T_0 = 2{,}7255$~K) avec l'incertitude $1\sigma$ propagée depuis $H_0$ et $\Omega_m$ \emph{corrélés}
($\rho = -0{,}976$ : l'oublier surestimerait l'incertitude) et, au choix, la comparaison avec SH0ES
($H_0 = 73{,}04$, tension de Hubble). Pour comparer : Planck 2015, WMAP 9 ans, un $\Lambda$CDM « de manuel »
(70~/ 0,3~/ 0,7 sans rayonnement) et un jeu personnalisé ($H_0$, $\Omega_m$, $\Omega_k$, rayonnement de Planck 2018).
Plages : $H_0$ de 40 à 100, $\Omega_m$ de 0,1 à 1, $\Omega_k$ de $-0{,}3$ à $0{,}3$ ($\pm 0{,}05$ pour Planck 2018) ;
une valeur hors plage est \emph{refusée avec une explication}, jamais corrigée en silence.

\paragraph{Saisie.} Un $z$ (virgule ou point), un des exemples (M\,87, 3C\,273, $z = 1$, $z = 2{,}34$, ULAS J1120,
GN-z11, réionisation, fond diffus), ou un nom d'objet : son redshift est demandé à SIMBAD (nom tolérant :
« m87 », « ngc4486 » ; plusieurs candidats : vous choisissez). $z \le 0$ (objet qui s'approche, comme M\,31) ne donne
pas de distance ; sous $z = 0{,}03$ la vitesse propre de l'objet domine (avertissement) ; au-delà de $z \approx 1090$
l'univers est opaque. Tableau et courbes s'exportent en CSV (« ; », UTF-8 ; valeurs brutes en Mpc, Gyr, Gpc$^3$).

\paragraph{Curseur.} Sous le champ $z$, un curseur parcourt tout le domaine sur une échelle logarithmique
(de 0,001 à 1100, repères 0,01 / 0,1 / 1 / 10 / 100 / 1000). En glissant, le marqueur vertical des courbes et le
tableau des valeurs suivent en direct : les valeurs sont alors \emph{interpolées} sur la grille des courbes (300
points logarithmiques, écart typique $< 0{,}5$~\%, colonnes $\sigma$ et SH0ES suspendues) et la ligne d'état le dit ;
au relâchement, le calcul exact reprend. Flèches : pas fin (1~\% d'une décade) ; Page haut/bas : une décade. Le
champ $z$ et le curseur restent synchronisés dans les deux sens.

\capture{img/module_cosmo_cote_fr.png}{1.0}
\paragraph{Disposition.} Côte à côte, le tableau prend la largeur de ses colonnes (au contenu, jamais tronquées) et
les courbes tout le reste ; empilés, le tableau montre ses 18 lignes si la place le permet, sinon la place est
partagée en gardant au moins 300~px aux courbes. \emph{Affichage > Disposition de la Cosmologie} (ou la liste
\emph{Disposition} du module) : automatique (côte à côte à partir de 1~550~px de large, retour à l'empilement sous
1~450~px : pas de bascule pendant un redimensionnement), toujours côte à côte, ou toujours empilés ; le choix et la
position du séparateur de chaque disposition sont gardés.

\paragraph{Unités.} Notation à exposants négatifs : $H_0$ en km\,s$^{-1}$\,Mpc$^{-1}$, vitesses en km\,s$^{-1}$ ;
dans les CSV, en ASCII non ambigu (\texttt{km s\^{}-1 Mpc\^{}-1}), de même sur une console qui ne sait pas afficher
les exposants (Windows cp1252).

\paragraph{Origine et vérification.} Le noyau de calcul est repris du calculateur « cosmologie-redshift » du même
auteur (dépôt public \url{https://github.com/ARP273-ROSE/cosmologie-redshift}), avec astropy.cosmology. Avant
l'intégration, il a été comparé le 8 octobre 2026 à astropy et à une intégration indépendante sous SageMath
(mpmath, 18 chiffres, fonction de Fermi-Dirac exacte pour les neutrinos), de $z = 10^{-8}$ à 1100, pour
$\Omega_k = 0, \pm0{,}01, \pm0{,}05$ : écart relatif $\le 3\times10^{-6}$ sur les distances et
$\le 2\times10^{-5}$ sur les âges et $E(z)$ (ajustement de Komatsu des neutrinos dans astropy) ; aucun bug. Le
volume comobile d'astropy perd en revanche tous ses chiffres à très petit $z$ en univers courbe (jusqu'à un volume
négatif) : Coupole emploie là un développement limité, raccordé à la forme fermée. Scripts :
\texttt{outils/verif\_cosmo/}.

\begin{ccode}
coupole cosmo 0,158 --shoes\\
coupole cosmo --objet "3C 273"\\
coupole cosmo 1 2 3 --modele wmap9 --csv resultats.csv\\
coupole cosmo --courbes 0,001 1500 300 --csv courbes.csv
\end{ccode}

% =====================================================================
\section{Machine et performance}
\capture{img/module_machine_fr.png}{0.9}
Détection sans dépendance obligatoire : système, architecture, processeur, cœurs physiques et logiques (bornés par
l'affinité et le quota du conteneur), mémoire totale et disponible, place libre, cartes graphiques (nvidia-smi,
\texttt{/sys/class/drm}, system\_profiler, CIM). Parallélisme : téléchargements 3 (4 au plus), conversions = cœurs
physiques − 1 bornées par la mémoire disponible à raison de 500 Mo par conversion (pic mesuré sur une image IRIS
4096², la moitié de la mémoire libre au plus) ; mode économe (1 conversion, 2 téléchargements) automatique sous 4 Go
ou 2 cœurs ; bridage manuel, lui aussi borné par la mémoire. Exemples : 2 cœurs et 2 Go → 2 téléchargements,
1 conversion ; 6 cœurs et 20 Go libres → 3 et 5 ; 32 cœurs et 60 Go libres → 3 et 16 (plafond). La carte
graphique n'est pas utilisée par la Banque OHP (travail limité par le réseau, le disque et Zstandard) ;
\texttt{coupole.core.calcul} offre CuPy aux modules futurs.

Rien de lourd ne s'exécute dans le fil graphique : astropy, scipy et SEP ne sont chargés qu'au premier besoin, les
sondes de la machine et la détection d'ASTAP tournent en fond, et la fermeture de la fenêtre arrête puis attend tous
les fils, tâches et processus (voir \texttt{docs/AUDIT\_2026-10.md}). Les grandes listes restent fluides à 80~000
lignes : les cellules ne sont calculées que pour les lignes affichées, le tri se fait sur une clé par ligne (et non par
comparaisons deux à deux), la sélection ne fait plus parcourir toutes les lignes à chaque dessin, et le changement de
thème ne recalcule que les couleurs (voir \texttt{docs/AUDIT2\_2026-10.md}).

\begin{center}\small
\begin{tabular}{@{}lr@{}}
\toprule
\textbf{Mesure (processeur à 6 cœurs, NVMe ; conteneur Debian)} & \textbf{Valeur}\\
\midrule
\texttt{coupole --version} & 0,1–0,2 s\\
fenêtre affichée (astropy, scipy, SEP non chargés) & 0,42 s (1,5 s avant l'audit)\\
inventaire chargé (fil de fond) & 3,3 s au premier lancement, 0,8 s ensuite\\
chaîne complète, 24 images T120 de 8 Mo, plafond 8 Mo/s & 24,5 s = durée du téléchargement seul (24,2 s)\\
même chaîne sans plafond (serveur local) & 2,3 s, 86 Mo/s de FITS, 10 images/s\\
conversion d'une image T120 (un processus) & 0,13 s (zstd 14 ms, lecture 18 ms, écriture et relecture 29 ms)\\
pic mémoire d'une conversion T120 / IRIS 4096² (processus dédié) & 163 Mo / 365 Mo (170 / 573 avant)\\
interface pendant 3 téléchargements (28 Mo/s) & écart max. du minuteur 20 ms : 20,8 ms\\
\midrule
\multicolumn{2}{@{}l}{\emph{Second audit (0.1.4) : banque réelle de 7~625 XISF ; partage réseau simulé à 2 ms par accès}}\\
rangement final sans rien à déplacer, sur le partage & 1,5 s (60,7 s avant)\\
48 images vers une destination sur le partage & 4,9 s (20,3 s avant)\\
qualité, relance de la banque déjà mesurée (partage) & 1,9 s (72 s avant)\\
réorganiser la banque (lecture des 7~625 en-têtes) & 57 s (243 s avant)\\
tout sélectionner, 80~000 images (10 × la banque) & 0,2 s (3,1 s avant ; 0,5 s → 0,17 s sur la banque)\\
trier 80~000 images & 0,1–0,2 s (3,7–4,1 s avant)\\
\bottomrule
\end{tabular}
\end{center}

% =====================================================================
\section{Ligne de commande}
Tout ce que fait l'interface se fait en ligne de commande ; \cmd{--json} pour les scripts. Sans argument,
\cmd{coupole} ouvre l'interface. Les commandes ont un nom français et un alias anglais.
\begin{ccode}
coupole ohp estimer "NGC 6888" --nuit 2025-07-16\\
coupole ohp traiter "NGC 6888" --nuit 2025-07-16 --filtre B --filtre V --filtre R --dest \textasciitilde/OHP --oui\\
coupole ohp telecharger Pluton --sans-dates-douteuses --dest \textasciitilde/FITS\\
coupole qualite \textasciitilde/OHP --ecrire
\end{ccode}
{\small% Généré par outils/tables_cli.py — ne pas modifier à la main.
\begin{tabularx}{\linewidth}{@{}>{\ttfamily\small\raggedright}p{0.38\linewidth}X@{}}
\toprule
\multicolumn{2}{@{}p{\linewidth}@{}}{\sffamily\bfseries coupole}\\
\midrule
--version & affiche la version et quitte\\
--lang,\allowbreak{} --langue \{auto,\allowbreak{}fr,\allowbreak{}en\} & langue des messages (auto : celle du système)\\
-v,\allowbreak{} --verbeux,\allowbreak{} --verbose & journal détaillé\\
--reinitialiser-interface,\allowbreak{} --reset-interface & efface la disposition gardée (fenêtre, colonnes, filtres, dossiers des dialogues) puis quitte ; les réglages ne changent pas\\
\bottomrule
\end{tabularx}
\par\medskip
\begin{tabularx}{\linewidth}{@{}>{\ttfamily\small\raggedright}p{0.38\linewidth}X@{}}
\toprule
\multicolumn{2}{@{}p{\linewidth}@{}}{\sffamily\bfseries coupole astap\newline\normalfont\small Cherche ASTAP et son catalogue d'étoiles, indique quoi installer pour ce système, ou enregistre leur emplacement. Coupole fonctionne sans ASTAP.}\\
\midrule
--guide & affiche le guide d'installation pour ce système\\
--definir,\allowbreak{} --set CHEMIN & chemin de l'exécutable ASTAP (ou de son dossier)\\
--catalogue,\allowbreak{} --catalog DOSSIER & dossier du catalogue d'étoiles (D80, D50...)\\
--oublier,\allowbreak{} --forget & oublie les chemins enregistrés (retour à la détection)\\
--json & sortie JSON (pour les scripts)\\
\bottomrule
\end{tabularx}
\par\medskip
\begin{tabularx}{\linewidth}{@{}>{\ttfamily\small\raggedright}p{0.38\linewidth}X@{}}
\toprule
\multicolumn{2}{@{}p{\linewidth}@{}}{\sffamily\bfseries coupole rapports (reports)\newline\normalfont\small Rien n'est envoyé sans accord. Sans accord, les rapports restent dans un dossier local.}\\
\midrule
--activer,\allowbreak{} --enable & autorise l'envoi des rapports (anonymes)\\
--desactiver,\allowbreak{} --disable & interdit l'envoi des rapports\\
\bottomrule
\end{tabularx}
\par\medskip
\begin{tabularx}{\linewidth}{@{}>{\ttfamily\small\raggedright}p{0.38\linewidth}X@{}}
\toprule
\multicolumn{2}{@{}p{\linewidth}@{}}{\sffamily\bfseries coupole maj (update)\newline\normalfont\small Interroge les versions publiées sur GitHub. Le paquet autonome se met à jour lui-même ; une installation par pip/pipx indique la commande à lancer.}\\
\midrule
--appliquer,\allowbreak{} --apply & télécharge et installe la mise à jour (paquet autonome)\\
\bottomrule
\end{tabularx}
\par\medskip
\begin{tabularx}{\linewidth}{@{}>{\ttfamily\small\raggedright}p{0.38\linewidth}X@{}}
\toprule
\multicolumn{2}{@{}p{\linewidth}@{}}{\sffamily\bfseries coupole sources\newline\normalfont\small Sans option : liste les adresses et leur origine (défaut livré, fichier distant, valeur forcée).}\\
\midrule
--forcer,\allowbreak{} --set CLÉ VALEUR & force une valeur (jamais écrasée par une mise à jour)\\
--defaut,\allowbreak{} --default CLÉ & revient à la valeur par défaut pour cette clé\\
--tout-reinitialiser,\allowbreak{} --reset-all & revient aux valeurs par défaut pour toutes les clés\\
--tester,\allowbreak{} --test CLÉ & teste la connexion à cette adresse\\
--distant,\allowbreak{} --remote & récupère maintenant le fichier de sources publié dans le dépôt\\
\bottomrule
\end{tabularx}
\par\medskip
\begin{tabularx}{\linewidth}{@{}>{\ttfamily\small\raggedright}p{0.38\linewidth}X@{}}
\toprule
\multicolumn{2}{@{}p{\linewidth}@{}}{\sffamily\bfseries coupole ohp inventaire (inventory)\newline\normalfont\small Inventaire de la banque par le service TAP public (http://tap-ufe.obspm.fr/tap, table ivoa.obscore), gardé en cache local.}\\
\midrule
--rafraichir,\allowbreak{} --refresh & réinterroge le service TAP de l'Observatoire (une requête)\\
--manquantes,\allowbreak{} --missing & liste ce qui reste à télécharger par rapport au dossier de sortie (--dest)\\
--dest DOSSIER & dossier de destination (défaut : \textasciitilde{}/Coupole/OHP\_DU\_ECU ou celui des réglages)\\
--json & sortie JSON (pour les scripts)\\
\bottomrule
\end{tabularx}
\par\medskip
\begin{tabularx}{\linewidth}{@{}>{\ttfamily\small\raggedright}p{0.38\linewidth}X@{}}
\toprule
\multicolumn{2}{@{}p{\linewidth}@{}}{\sffamily\bfseries coupole ohp catalogue (catalog)\newline\normalfont\small Les 194 noms de la base regroupés en objets, classés par type.}\\
\midrule
--type TYPE & catégorie : ast, neocp, com, pla, tno, pn, neb, amas, gal, eto (répétable)\\
--telescope \{T120,\allowbreak{}IRIS\} & télescope (répétable)\\
--chercher,\allowbreak{} --search TEXTE & filtre les objets dont le nom contient ce texte\\
--json & sortie JSON (pour les scripts)\\
\bottomrule
\end{tabularx}
\par\medskip
\begin{tabularx}{\linewidth}{@{}>{\ttfamily\small\raggedright}p{0.38\linewidth}X@{}}
\toprule
\multicolumn{2}{@{}p{\linewidth}@{}}{\sffamily\bfseries coupole ohp images (list)\newline\normalfont\small Une ligne par image : date, télescope, filtre, pose, drapeaux (doublon, date partagée, plein jour).}\\
\midrule
OBJET & objet(s) : « NGC 6888 », « Pluton », « 1998 AG6 », ou un nom de la base\\
--type TYPE & catégorie : ast, neocp, com, pla, tno, pn, neb, amas, gal, eto (répétable)\\
--telescope \{T120,\allowbreak{}IRIS\} & télescope (répétable)\\
--filtre,\allowbreak{} --filter F & filtre : B, V, R (Johnson et Cousins), Bc, Vc, Rc, Ha, OIII, g, r, i, z, ou le nom exact (« Johnson R ») ; répétable ou séparé par des virgules\\
--nuit,\allowbreak{} --night AAAA-MM-JJ & nuit (date du soir) ; répétable\\
--sans-dates-douteuses,\allowbreak{} --no-doubtful-dates & écarte les poses à DATE-OBS partagée ou datées en plein jour\\
--tout,\allowbreak{} --all & toute la banque (≈ 78 Go)\\
--csv FICHIER & écrit aussi la liste dans ce fichier CSV\\
--dest DOSSIER & dossier de destination (défaut : \textasciitilde{}/Coupole/OHP\_DU\_ECU ou celui des réglages)\\
\bottomrule
\end{tabularx}
\par\medskip
\begin{tabularx}{\linewidth}{@{}>{\ttfamily\small\raggedright}p{0.38\linewidth}X@{}}
\toprule
\multicolumn{2}{@{}p{\linewidth}@{}}{\sffamily\bfseries coupole ohp estimer (estimate)\newline\normalfont\small Estimation avant téléchargement, avec la place libre à destination.}\\
\midrule
OBJET & objet(s) : « NGC 6888 », « Pluton », « 1998 AG6 », ou un nom de la base\\
--type TYPE & catégorie : ast, neocp, com, pla, tno, pn, neb, amas, gal, eto (répétable)\\
--telescope \{T120,\allowbreak{}IRIS\} & télescope (répétable)\\
--filtre,\allowbreak{} --filter F & filtre : B, V, R (Johnson et Cousins), Bc, Vc, Rc, Ha, OIII, g, r, i, z, ou le nom exact (« Johnson R ») ; répétable ou séparé par des virgules\\
--nuit,\allowbreak{} --night AAAA-MM-JJ & nuit (date du soir) ; répétable\\
--sans-dates-douteuses,\allowbreak{} --no-doubtful-dates & écarte les poses à DATE-OBS partagée ou datées en plein jour\\
--tout,\allowbreak{} --all & toute la banque (≈ 78 Go)\\
--format \{xisf,\allowbreak{}xisf16,\allowbreak{}fz,\allowbreak{}fits\} & format de sortie : xisf (PixInsight, défaut), xisf16 (XISF compatible N.I.N.A. et Siril : entiers 16 bits, zlib), fz (FITS compressé sans perte), fits (float32)\\
--dest DOSSIER & dossier de destination (défaut : \textasciitilde{}/Coupole/OHP\_DU\_ECU ou celui des réglages)\\
\bottomrule
\end{tabularx}
\par\medskip
\begin{tabularx}{\linewidth}{@{}>{\ttfamily\small\raggedright}p{0.38\linewidth}X@{}}
\toprule
\multicolumn{2}{@{}p{\linewidth}@{}}{\sffamily\bfseries coupole ohp traiter (process)\newline\normalfont\small Traitement complet d'une sélection, reprenable : relancer la même commande reprend là où elle s'est arrêtée.}\\
\midrule
OBJET & objet(s) : « NGC 6888 », « Pluton », « 1998 AG6 », ou un nom de la base\\
--type TYPE & catégorie : ast, neocp, com, pla, tno, pn, neb, amas, gal, eto (répétable)\\
--telescope \{T120,\allowbreak{}IRIS\} & télescope (répétable)\\
--filtre,\allowbreak{} --filter F & filtre : B, V, R (Johnson et Cousins), Bc, Vc, Rc, Ha, OIII, g, r, i, z, ou le nom exact (« Johnson R ») ; répétable ou séparé par des virgules\\
--nuit,\allowbreak{} --night AAAA-MM-JJ & nuit (date du soir) ; répétable\\
--sans-dates-douteuses,\allowbreak{} --no-doubtful-dates & écarte les poses à DATE-OBS partagée ou datées en plein jour\\
--tout,\allowbreak{} --all & toute la banque (≈ 78 Go)\\
--dest DOSSIER & dossier de destination (défaut : \textasciitilde{}/Coupole/OHP\_DU\_ECU ou celui des réglages)\\
--format \{xisf,\allowbreak{}xisf16,\allowbreak{}fz,\allowbreak{}fits\} & format de sortie : xisf (PixInsight, défaut), xisf16 (XISF compatible N.I.N.A. et Siril : entiers 16 bits, zlib), fz (FITS compressé sans perte), fits (float32)\\
--noms,\allowbreak{} --names \{fr,\allowbreak{}en\} & langue des noms de dossiers, d'objets et des en-têtes (défaut : langue de l'interface)\\
--astap \{auto,\allowbreak{}tous,\allowbreak{}suspectes,\allowbreak{}jamais,\allowbreak{}all,\allowbreak{}suspicious,\allowbreak{}never\} & vérification ASTAP : auto/tous (chaque image si ASTAP est là), suspectes, jamais\\
--telechargements,\allowbreak{} --downloads N & téléchargements simultanés (défaut automatique ; 4 au plus)\\
--conversions N & conversions simultanées (défaut : selon cœurs et mémoire)\\
--econome,\allowbreak{} --eco & mode économe : 1 conversion, 2 téléchargements\\
--debit,\allowbreak{} --rate Mo/s & débit maximal vers le serveur (défaut 8 Mo/s)\\
--garder-fits,\allowbreak{} --keep-fits & garde aussi les FITS d'origine (dans \_traitement/fits\_origine)\\
--garder-doublons,\allowbreak{} --keep-duplicates & traite aussi les doublons (inventaire et pixels identiques) au lieu de les écarter\\
--oui,\allowbreak{} --yes & accepte une sélection de plus de 5 Go sans confirmation\\
\bottomrule
\end{tabularx}
\par\medskip
\begin{tabularx}{\linewidth}{@{}>{\ttfamily\small\raggedright}p{0.38\linewidth}X@{}}
\toprule
\multicolumn{2}{@{}p{\linewidth}@{}}{\sffamily\bfseries coupole ohp tout (all)\newline\normalfont\small Toute la banque (≈ 8 000 images, ≈ 78 Go de FITS → ≈ 30 Go en XISF), avec estimation du volume et du temps avant confirmation ; relancer la même commande reprend là où elle s'était arrêtée.}\\
\midrule
--dest DOSSIER & dossier de destination (défaut : \textasciitilde{}/Coupole/OHP\_DU\_ECU ou celui des réglages)\\
--format \{xisf,\allowbreak{}xisf16,\allowbreak{}fz,\allowbreak{}fits\} & format de sortie : xisf (PixInsight, défaut), xisf16 (XISF compatible N.I.N.A. et Siril : entiers 16 bits, zlib), fz (FITS compressé sans perte), fits (float32)\\
--noms,\allowbreak{} --names \{fr,\allowbreak{}en\} & langue des noms de dossiers, d'objets et des en-têtes (défaut : langue de l'interface)\\
--astap \{auto,\allowbreak{}tous,\allowbreak{}suspectes,\allowbreak{}jamais,\allowbreak{}all,\allowbreak{}suspicious,\allowbreak{}never\} & vérification ASTAP : auto/tous (chaque image si ASTAP est là), suspectes, jamais\\
--telechargements,\allowbreak{} --downloads N & téléchargements simultanés (défaut automatique ; 4 au plus)\\
--conversions N & conversions simultanées (défaut : selon cœurs et mémoire)\\
--econome,\allowbreak{} --eco & mode économe : 1 conversion, 2 téléchargements\\
--debit,\allowbreak{} --rate Mo/s & débit maximal vers le serveur (défaut 8 Mo/s)\\
--garder-fits,\allowbreak{} --keep-fits & garde aussi les FITS d'origine (dans \_traitement/fits\_origine)\\
--garder-doublons,\allowbreak{} --keep-duplicates & traite aussi les doublons (inventaire et pixels identiques) au lieu de les écarter\\
--oui,\allowbreak{} --yes & accepte une sélection de plus de 5 Go sans confirmation\\
\bottomrule
\end{tabularx}
\par\medskip
\begin{tabularx}{\linewidth}{@{}>{\ttfamily\small\raggedright}p{0.38\linewidth}X@{}}
\toprule
\multicolumn{2}{@{}p{\linewidth}@{}}{\sffamily\bfseries coupole ohp nouveautes (new)\newline\normalfont\small Réinterroge le service TAP, compare à la copie locale (dossier de sortie) et liste les nouveautés ; --telecharger les traite dans la même arborescence (jamais sans cette option).}\\
\midrule
--telecharger,\allowbreak{} --download & télécharge et convertit les nouveautés listées\\
--dest DOSSIER & dossier de destination (défaut : \textasciitilde{}/Coupole/OHP\_DU\_ECU ou celui des réglages)\\
--format \{xisf,\allowbreak{}xisf16,\allowbreak{}fz,\allowbreak{}fits\} & format de sortie : xisf (PixInsight, défaut), xisf16 (XISF compatible N.I.N.A. et Siril : entiers 16 bits, zlib), fz (FITS compressé sans perte), fits (float32)\\
--noms,\allowbreak{} --names \{fr,\allowbreak{}en\} & langue des noms de dossiers, d'objets et des en-têtes (défaut : langue de l'interface)\\
--astap \{auto,\allowbreak{}tous,\allowbreak{}suspectes,\allowbreak{}jamais,\allowbreak{}all,\allowbreak{}suspicious,\allowbreak{}never\} & vérification ASTAP : auto/tous (chaque image si ASTAP est là), suspectes, jamais\\
--telechargements,\allowbreak{} --downloads N & téléchargements simultanés (défaut automatique ; 4 au plus)\\
--conversions N & conversions simultanées (défaut : selon cœurs et mémoire)\\
--econome,\allowbreak{} --eco & mode économe : 1 conversion, 2 téléchargements\\
--debit,\allowbreak{} --rate Mo/s & débit maximal vers le serveur (défaut 8 Mo/s)\\
--garder-fits,\allowbreak{} --keep-fits & garde aussi les FITS d'origine (dans \_traitement/fits\_origine)\\
--garder-doublons,\allowbreak{} --keep-duplicates & traite aussi les doublons (inventaire et pixels identiques) au lieu de les écarter\\
--oui,\allowbreak{} --yes & accepte une sélection de plus de 5 Go sans confirmation\\
\bottomrule
\end{tabularx}
\par\medskip
\begin{tabularx}{\linewidth}{@{}>{\ttfamily\small\raggedright}p{0.38\linewidth}X@{}}
\toprule
\multicolumn{2}{@{}p{\linewidth}@{}}{\sffamily\bfseries coupole ohp reorganiser (reorganise, reorganize)\newline\normalfont\small Parcourt un dossier, reconnaît les fichiers produits par Coupole (en-tête) et les déplace dans l'arborescence des lots de --dest ; conflit de nom → suffixe et ligne de journal ; rien n'est écrasé.}\\
\midrule
DOSSIER & dossier où se trouvent les fichiers à ranger\\
--dest DOSSIER & dossier de destination (défaut : \textasciitilde{}/Coupole/OHP\_DU\_ECU ou celui des réglages)\\
\bottomrule
\end{tabularx}
\par\medskip
\begin{tabularx}{\linewidth}{@{}>{\ttfamily\small\raggedright}p{0.38\linewidth}X@{}}
\toprule
\multicolumn{2}{@{}p{\linewidth}@{}}{\sffamily\bfseries coupole ohp telecharger (download)\newline\normalfont\small Rangement : <dossier>/<objet>/<nuit>\_<télescope>/<fichier d'origine>.}\\
\midrule
OBJET & objet(s) : « NGC 6888 », « Pluton », « 1998 AG6 », ou un nom de la base\\
--type TYPE & catégorie : ast, neocp, com, pla, tno, pn, neb, amas, gal, eto (répétable)\\
--telescope \{T120,\allowbreak{}IRIS\} & télescope (répétable)\\
--filtre,\allowbreak{} --filter F & filtre : B, V, R (Johnson et Cousins), Bc, Vc, Rc, Ha, OIII, g, r, i, z, ou le nom exact (« Johnson R ») ; répétable ou séparé par des virgules\\
--nuit,\allowbreak{} --night AAAA-MM-JJ & nuit (date du soir) ; répétable\\
--sans-dates-douteuses,\allowbreak{} --no-doubtful-dates & écarte les poses à DATE-OBS partagée ou datées en plein jour\\
--tout,\allowbreak{} --all & toute la banque (≈ 78 Go)\\
--dest DOSSIER & dossier de destination (défaut : \textasciitilde{}/Coupole/OHP\_DU\_ECU ou celui des réglages)\\
--telechargements,\allowbreak{} --downloads N & téléchargements simultanés (défaut automatique ; 4 au plus)\\
--debit,\allowbreak{} --rate Mo/s & débit maximal vers le serveur (défaut 8 Mo/s)\\
--oui,\allowbreak{} --yes & accepte une sélection de plus de 5 Go sans confirmation\\
\bottomrule
\end{tabularx}
\par\medskip
\begin{tabularx}{\linewidth}{@{}>{\ttfamily\small\raggedright}p{0.38\linewidth}X@{}}
\toprule
\multicolumn{2}{@{}p{\linewidth}@{}}{\sffamily\bfseries coupole ohp anomalies (anomalies-report)\newline\normalfont\small Doublons (même fichier, copie diurne, pixels identiques), dates partagées ou diurnes, champ incohérent, classement à vérifier, instrument inconnu. Rien n'est supprimé.}\\
\midrule
--dest DOSSIER & ajoute les doublons de pixels trouvés dans ce dossier de traitement\\
--csv FICHIER & écrit aussi la liste dans ce fichier CSV\\
--nouveaux,\allowbreak{} --new & seulement les images nouvelles depuis le dernier rafraîchissement\\
\bottomrule
\end{tabularx}
\par\medskip
\begin{tabularx}{\linewidth}{@{}>{\ttfamily\small\raggedright}p{0.38\linewidth}X@{}}
\toprule
\multicolumn{2}{@{}p{\linewidth}@{}}{\sffamily\bfseries coupole ohp classer (classify)\newline\normalfont\small Sans argument : liste les noms à vérifier. Avec un nom de la base et --objet : mémorise la correction (fusion avec un objet existant si --type est omis). --en-ligne : interroge JPL SBDB et CDS Sesame.}\\
\midrule
NOM\_BASE & nom de cible tel qu'il figure dans la base\\
--objet,\allowbreak{} --object OBJET & objet canonique à associer (existant : fusion)\\
--type \{ast,\allowbreak{}neocp,\allowbreak{}com,\allowbreak{}pla,\allowbreak{}tno,\allowbreak{}pn,\allowbreak{}neb,\allowbreak{}amas,\allowbreak{}gal,\allowbreak{}eto,\allowbreak{}autre\} & catégorie : ast, neocp, com, pla, tno, pn, neb, amas, gal, eto (répétable)\\
--oublier,\allowbreak{} --forget & oublie la correction de ce nom\\
--en-ligne,\allowbreak{} --online & résout en ligne les noms inconnus (JPL SBDB, CDS Sesame), avec cache\\
\bottomrule
\end{tabularx}
\par\medskip
\begin{tabularx}{\linewidth}{@{}>{\ttfamily\small\raggedright}p{0.38\linewidth}X@{}}
\toprule
\multicolumn{2}{@{}p{\linewidth}@{}}{\sffamily\bfseries coupole ohp ranger (sort)\newline\normalfont\small Utile après une interruption, ou pour regrouper plusieurs traitements dans un même dossier.}\\
\midrule
--dest DOSSIER & dossier de destination (défaut : \textasciitilde{}/Coupole/OHP\_DU\_ECU ou celui des réglages)\\
\bottomrule
\end{tabularx}
\par\medskip
\begin{tabularx}{\linewidth}{@{}>{\ttfamily\small\raggedright}p{0.38\linewidth}X@{}}
\toprule
\multicolumn{2}{@{}p{\linewidth}@{}}{\sffamily\bfseries coupole ohp bilan (status)\newline\normalfont\small état d'avancement d'un dossier de traitement}\\
\midrule
--dest DOSSIER & dossier de destination (défaut : \textasciitilde{}/Coupole/OHP\_DU\_ECU ou celui des réglages)\\
--json & sortie JSON (pour les scripts)\\
\bottomrule
\end{tabularx}
\par\medskip
\begin{tabularx}{\linewidth}{@{}>{\ttfamily\small\raggedright}p{0.38\linewidth}X@{}}
\toprule
\multicolumn{2}{@{}p{\linewidth}@{}}{\sffamily\bfseries coupole ohp metadonnees (metadata)\newline\normalfont\small Pour chaque XISF/FITS du dossier : FOCALLEN accordé à l'échelle mesurée (ancienne valeur en HISTORY) et propriétés XISF Instrument:Sensor:XPixelSize, Instrument:Camera:XBinning, Instrument:Telescope:FocalLength… Les pixels ne sont pas touchés (relus et comparés), écriture atomique, journal dans \_traitement/metadonnees.csv. Sans --reecrire : simulation.}\\
\midrule
DOSSIER & dossier de sortie (ou un lot)\\
--reecrire,\allowbreak{} --rewrite & écrit vraiment (sinon : dit seulement ce qui changerait)\\
--noms,\allowbreak{} --names \{fr,\allowbreak{}en\} & langue des noms de dossiers, d'objets et des en-têtes (défaut : langue de l'interface)\\
\bottomrule
\end{tabularx}
\par\medskip
\begin{tabularx}{\linewidth}{@{}>{\ttfamily\small\raggedright}p{0.38\linewidth}X@{}}
\toprule
\multicolumn{2}{@{}p{\linewidth}@{}}{\sffamily\bfseries coupole qualite\newline\normalfont\small Facultatif : fond, bruit, FWHM et ellipticité (carte 3×3), gradient du fond, saturation, traînées, échantillonnage}\\
\midrule
DOSSIER & dossier de lots (ou un fichier image)\\
--ecrire,\allowbreak{} --write & écrit QUALITE.csv et QUALITE.txt dans chaque lot\\
--echantillon,\allowbreak{} --sample N & mesure N images par lot (réparties dans le temps) au lieu de toutes\\
--tout,\allowbreak{} --all & tout mesurer même au-delà de 200 images (sinon : échantillon de 5 par lot, annoncé)\\
--processus,\allowbreak{} --processes N & nombre de processus de mesure (défaut : plan machine)\\
\bottomrule
\end{tabularx}
\par\medskip
\begin{tabularx}{\linewidth}{@{}>{\ttfamily\small\raggedright}p{0.38\linewidth}X@{}}
\toprule
\multicolumn{2}{@{}p{\linewidth}@{}}{\sffamily\bfseries coupole donnees lire (read)\newline\normalfont\small décrit le contenu d'un fichier (FITS, CSV)}\\
\midrule
FICHIER & fichier à lire\\
--json & sortie JSON (pour les scripts)\\
\bottomrule
\end{tabularx}
\par\medskip
\begin{tabularx}{\linewidth}{@{}>{\ttfamily\small\raggedright}p{0.38\linewidth}X@{}}
\toprule
\multicolumn{2}{@{}p{\linewidth}@{}}{\sffamily\bfseries coupole donnees exporter (export)\newline\normalfont\small exporte un spectre ou une série en CSV}\\
\midrule
FICHIER & fichier à lire\\
--csv FICHIER & exporte un spectre ou une série en CSV\\
--hdu HDU & numéro de la donnée dans le fichier (défaut : la première 1D)\\
--vitesse,\allowbreak{} --velocity & ajoute la vitesse radio (spectre en fréquence)\\
--f0 MHz & fréquence de repos en MHz (défaut : RESTFRQ, sinon H I)\\
\bottomrule
\end{tabularx}
\par\medskip
\begin{tabularx}{\linewidth}{@{}>{\ttfamily\small\raggedright}p{0.38\linewidth}X@{}}
\toprule
\multicolumn{2}{@{}p{\linewidth}@{}}{\sffamily\bfseries coupole cosmo\newline\normalfont\small Du redshift aux distances : comobile, luminosité, angulaire, temps de regard en arrière, âge, volume, module de distance, échelle (Planck 2018 et autres jeux de paramètres)}\\
\midrule
Z & redshift(s) à calculer\\
--objet,\allowbreak{} --object NOM & nom d'objet : redshift demandé à SIMBAD\\
--modele,\allowbreak{} --model \{planck18,\allowbreak{}planck15,\allowbreak{}wmap9,\allowbreak{}simple,\allowbreak{}perso\} & jeu de paramètres (défaut : planck18)\\
--h0 H0 & H₀ en km s⁻¹ Mpc⁻¹ (modèle perso)\\
--om OM & Ωm total (modèle perso)\\
--ok OK & Ωk, courbure (planck18 ou perso)\\
--shoes & ajoute la comparaison SH0ES (planck18)\\
--sans-incertitudes,\allowbreak{} --no-uncertainties & sans les incertitudes (plus rapide)\\
--csv FICHIER & écrit les résultats dans ce fichier CSV\\
--courbes,\allowbreak{} --curves ZMIN ZMAX N & écrit les courbes (z de ZMIN à ZMAX, N points) dans le fichier CSV donné par --csv\\
--json & sortie JSON (pour les scripts)\\
\bottomrule
\end{tabularx}
\par\medskip
\begin{tabularx}{\linewidth}{@{}>{\ttfamily\small\raggedright}p{0.38\linewidth}X@{}}
\toprule
\multicolumn{2}{@{}p{\linewidth}@{}}{\sffamily\bfseries coupole sites\newline\normalfont\small Sites d'observation (carte OpenStreetMap), fuseaux horaires, heure UTC et heure locale}\\
\midrule
--ajouter,\allowbreak{} --add X & ajoute un site : ID NOM LAT LON ALT FUSEAU [MPC]\\
--supprimer,\allowbreak{} --delete ID & supprime un site ajouté\\
--heure,\allowbreak{} --time DATE\_ISO & affiche une date UTC (ISO) en heure du site et date du soir\\
--site ID & identifiant du site (défaut : ohp)\\
\bottomrule
\end{tabularx}
\par\medskip
\begin{tabularx}{\linewidth}{@{}>{\ttfamily\small\raggedright}p{0.38\linewidth}X@{}}
\toprule
\multicolumn{2}{@{}p{\linewidth}@{}}{\sffamily\bfseries coupole machine\newline\normalfont\small Diagnostic matériel, parallélisme retenu, carte graphique, ASTAP}\\
\midrule
--json & sortie JSON (pour les scripts)\\
\bottomrule
\end{tabularx}
\par\medskip
}

% =====================================================================
\section{Évolutivité}\label{chap:evolutif}
\paragraph{Sources.} Les adresses (TAP, collection, portail, réécriture des URL de téléchargement, résolveurs,
point de collecte, dépôt, tuiles de carte) sont dans \texttt{coupole/donnees/sources.json}. Trois couches : valeurs
livrées ; fichier publié dans le dépôt, récupéré au démarrage s'il est plus récent et passe le contrôle de forme
(clés connues, types, domaines d'une liste blanche ; pas de signature cryptographique) ; valeurs forcées par
l'utilisateur, jamais écrasées. Préférences > Sources : modifier, revenir au défaut, tester. Si le serveur des
images déménage, \texttt{ohp.telechargement.de}/\texttt{vers} réécrit les URL.
\paragraph{Classement.} Table livrée \texttt{data/classement.json} (194 noms vérifiés) ; alias publiés ; corrections
de l'utilisateur (\emph{Corriger le classement}, \cmd{coupole ohp classer}) ; règles sur la forme du nom ; résolution
en ligne facultative (JPL SBDB, CDS Sesame), en cache ; regroupement par position sur un objet fixe connu ; sinon
« non classé ». Tout classement automatique est marqué « à vérifier ».
\paragraph{Base.} Colonnes ObsCore en plus ignorées, en moins remplacées par une valeur par défaut, instruments et
filtres inconnus rangés à part, données non images jamais signalées « de jour » ; testé sur inventaires simulés.
\paragraph{Sécurité.} Aucun secret dans le dépôt ; fichier de sources distant et point de collecte en HTTPS
obligatoire, taille bornée, liste blanche de domaines ; archive de mise à jour téléchargée seulement depuis GitHub,
contrôlée avant extraction (chemins absolus ou « .. » refusés, liens symboliques refusés, taille dépliée et nombre de
membres bornés), remplacement avec retour arrière en cas d'échec ; XISF et FITS lus avec des bornes de taille ; noms
de dossiers construits depuis la base rendus sûrs (ASCII, caractères interdits, noms réservés de Windows, longueur) ;
ASTAP lancé sans interpréteur de commandes ; aucun \texttt{eval} ni \texttt{pickle} sur des données externes.

% =====================================================================
\section{Rapports d'incident et confidentialité}
Au premier lancement, Coupole demande s'il peut envoyer des rapports anonymes : installation, preuve de vie (une
fois par jour), plantage, crash natif, interface figée, signalement écrit. Contenu : version, système, processeur,
mémoire, carte graphique, identifiant d'installation tiré au hasard, trace d'erreur avec le dossier personnel
remplacé par « \textasciitilde », le nom d'utilisateur et le nom de la machine masqués et les chemins profonds
tronqués ; un rapport ne dépasse jamais 64 ko. Jamais de nom de machine, d'utilisateur, ni d'image. Sans accord, rien ne part : les rapports restent dans \texttt{\_rapports} du
dossier de réglages. Choix modifiable (menu Aide, \cmd{coupole rapports --activer|--desactiver}).

% =====================================================================
\section{Écrire un module}
Un module est un paquet Python avec un \texttt{manifest.json} (identifiant, version, nom en français et en anglais,
icône, textes, points d'entrée graphique et ligne de commande), découvert dans \texttt{coupole/modules/}, par point
d'entrée \texttt{coupole.modules}, ou dans le dossier \texttt{modules/} des réglages. Textes par \texttt{tr()},
info-bulles obligatoires, travail long hors du fil graphique : voir CONTRIBUTING.md.

% =====================================================================
\section{Dépannage}
\begin{center}\small
\begin{tabularx}{\linewidth}{@{}>{\raggedright}p{0.34\linewidth}X@{}}
\toprule
\textbf{Symptôme} & \textbf{Que faire}\\
\midrule
Rien ne démarre (Windows) & menu Démarrer > Coupole > Diagnostic : l'erreur reste affichée.\\
« Interface graphique indisponible » & bibliothèques Qt manquantes (Linux : \texttt{libxcb-cursor0}, \texttt{libegl1}) ; la ligne de commande reste utilisable.\\
Téléchargement en échec & relancer (reprise) ; vérifier la connexion ; Préférences > Sources > Tester ; serveur déménagé : réécriture d'URL.\\
« ASTAP non trouvé » & assistant ASTAP ; ou \texttt{coupole astap --definir CHEMIN --catalogue DOSSIER}.\\
Catalogue incomplet & réinstaller le catalogue (1~476 tuiles attendues).\\
Place insuffisante & choisir un autre dossier de sortie ; la place nécessaire inclut les FITS en attente.\\
Machine lente & mode économe, ou moins de conversions simultanées.\\
Nom de cible inconnu & onglet Catalogue, case « À vérifier », \emph{Corriger le classement}.\\
Qualité indisponible & \texttt{pip install sep} (inclus dans les paquets).\\
Réglages revenus par défaut & le fichier \texttt{reglages.json} était illisible : il est gardé à côté avec le suffixe \texttt{.corrompu-…}.\\
Dossier de sortie non inscriptible & message clair au lancement ; choisir un autre dossier (Préférences ou champ Dossier).\\
Dialogue de fichiers sans le NAS, ou générique & Linux : installer \texttt{xdg-desktop-portal} et le moteur du bureau (\texttt{-kde}, \texttt{-gnome} ou \texttt{-gtk}), relancer ; KDE : \texttt{kio-fuse} ; sinon taper le chemin du partage monté dans le champ (section~\ref{sec:dialogues}).\\
Dialogue du système qui ne s'ouvre pas & 0.1.10 et avant, Spectres et séries sous KDE : mettre à jour (filtre sans nom refusé par le portail) ; sinon Préférences > Boîtes de dialogue de fichiers : Qt ; chaque ouverture est notée dans \texttt{coupole.log}.\\
Images possédées mais rien à ouvrir (copie d'un ancien traitement) & mettre à jour (0.1.11) : les emplacements sont retrouvés par \texttt{journal.csv} (section~\ref{sec:ouvrir}).\\
Qualité : colonnes vides sur ses propres poses & mettre à jour (0.1.11) : lecteur XISF tolérant ; la colonne \emph{état} donne le motif.\\
\texttt{\textasciitilde/.local/bin/coupole} introuvable & ajouter \texttt{\textasciitilde/.local/bin} au \texttt{PATH} (ligne donnée par l'installeur).\\
Fichier \texttt{.part} ou \texttt{.tmp} dans \texttt{\_traitement} & reste d'une coupure : la reprise le réutilise (\texttt{.part}) ou le remplace ; rien à faire.\\
WBPP ne trouve pas focale et pixel & saisir les valeurs de l'encadré \emph{Pour PixInsight / N.I.N.A.} (onglet Lots) dans le panneau \emph{Astrometric solution} ; fichiers anciens : \cmd{coupole ohp metadonnees DOSSIER --reecrire}.\\
Image blanche dans N.I.N.A. ou Siril & format XISF compatible (onglet Traitement), section~\ref{sec:compat}.\\
« database is locked » vers un NAS (0.1.9 et avant) & mettre à jour : la base de travail locale règle le problème (section~\ref{sec:partage}).\\
« Deux versions de la base d'état » & deux ordinateurs (ou le NAS) ont traité vers le même dossier : accepter la fusion, ou \cmd{coupole ohp fusionner --dest DOSSIER} ; les deux bases d'origine sont gardées.\\
« Recopie de la base d'état sur le partage impossible » & partage débranché ou base ouverte ailleurs : rien n'est perdu, nouvel essai au passage suivant et au prochain lancement.\\
\bottomrule
\end{tabularx}
\end{center}

\section*{Sources}
\addcontentsline{toc}{section}{Sources}
Banque « OHP student observations », Observatoire de Paris / PADC (Paris Astronomical Data Centre), service TAP
public, licence Etalab 2.0 ; guide officiel de la base (PADC-BDD-DU-ECU). Spécification XISF 1.0 (Pleiades
Astrophoto). ASTAP : \url{https://www.hnsky.org/astap.htm}. Codes d'observatoire : Minor Planet Center. SEP :
Barbary (2016), Bertin \& Arnouts (1996). Moments adaptatifs : Bernstein \& Jarvis (2002). Fond de carte :
© OpenStreetMap contributors. SIMBAD et Sesame : CDS, Strasbourg (Wenger et al. 2000) ; Aladin Lite (CDS) ;
NED (NASA/IPAC) ; JPL Small-Body Database (NASA/JPL). Cosmologie : Planck Collaboration (2020, A\&A 641, A6) ;
Hinshaw et al. (2013, WMAP 9) ; Riess et al. (2022, SH0ES) ; Hogg (1999, astro-ph/9905116) ; Komatsu et al. (2011)
pour la densité des neutrinos massifs ; astropy.cosmology ; calculateur « cosmologie-redshift » (même auteur).
\end{document}
